\documentclass[11pt]{article}
\usepackage[a4paper,top=2.4cm,bottom=2.5cm,left=2.0cm,right=2.0cm]{geometry}
\usepackage{amsmath,amssymb,amsfonts}
\usepackage{fontspec}
\setmainfont{TeX Gyre Termes}
\usepackage{booktabs}
\usepackage{graphicx}
\usepackage{pdflscape}
\graphicspath{{./}}
\usepackage[section]{placeins}
\usepackage{float}
\renewcommand{\topfraction}{0.92}
\renewcommand{\bottomfraction}{0.8}
\renewcommand{\textfraction}{0.07}
\renewcommand{\floatpagefraction}{0.75}
\usepackage{xcolor}
\definecolor{revisionblue}{RGB}{0,76,153}
\newcommand{\revmark}[1]{\textcolor{revisionblue}{\textbf{[R#1]}}\ }
\newcommand{\revdeleted}[1]{\par\noindent\textcolor{revisionblue}{\small\itshape [#1]}\par}
\newcommand{\revblue}[1]{\textcolor{revisionblue}{#1}}
\usepackage{tikz}
\usetikzlibrary{arrows.meta,positioning}
\usepackage{caption}
\captionsetup{font=small,labelfont=bf}
\usepackage{titlesec}
\titleformat{\section}{\normalfont\large\bfseries}{\thesection}{0.6em}{}
\titleformat{\subsection}{\normalfont\normalsize\bfseries}{\thesubsection}{0.5em}{}
\titleformat{\subsubsection}{\normalfont\normalsize\bfseries}{\thesubsubsection}{0.5em}{}
\titlespacing*{\section}{0pt}{1.3ex plus .2ex minus .2ex}{0.6ex}
\titlespacing*{\subsection}{0pt}{1.0ex plus .2ex minus .2ex}{0.5ex}
\usepackage{fancyhdr}
\pagestyle{fancy}
\fancyhf{}
\fancyhead[L]{\small\itshape Experimental Astronomy}
\fancyhead[R]{\small\thepage}
\renewcommand{\headrulewidth}{0.4pt}
\usepackage{cite}
\usepackage{hyperref}
\hypersetup{hidelinks}
\setlength{\emergencystretch}{3em}
\newcommand{\keV}{\ensuremath{\,\mathrm{keV}}}
\newcommand{\MeV}{\ensuremath{\,\mathrm{MeV}}}
\newcommand{\cms}{\ensuremath{\,\mathrm{cm^{-2}\,s^{-1}}}}
\newcommand{\phcms}{\ensuremath{\,\mathrm{ph\,cm^{-2}\,s^{-1}}}}
\newcommand{\cps}{\ensuremath{\,\mathrm{s^{-1}}}}
\newcommand{\wii}{W_{511}}
\newcommand{\aeff}{A_{\mathrm{eff}}}
\newcommand{\fref}{F_{\mathrm{ref}}}
\newcommand{\fthree}{F_{3\sigma}}
\raggedbottom

\begin{document}
\thispagestyle{fancy}
\begin{flushleft}
{\LARGE\bfseries Detector-coupled Monte Carlo estimate of background and unresolved-line sensitivity for a balloon-borne 511 keV Laue-lens TES telescope\par}
\vspace{1.1em}
{\large Authors omitted for review\par}
\vspace{0.3em}
{\normalsize\itshape Affiliations omitted for review\par}
\end{flushleft}
\vspace{0.7em}
\noindent\rule{\textwidth}{0.4pt}
\vspace{0.5em}

\noindent\textbf{Abstract}\hspace{0.7em}
Focused observations of compact or unresolved Galactic 511 keV emission with a
balloon-borne Laue lens require detector backgrounds to be controlled within a
sub-keV energy window. We developed an end-to-end Geant4/MEGAlib simulation in
which the focused signal, prompt atmospheric backgrounds, and activation decays
share a common detector-response and event-selection framework for a
transition-edge sensor (TES) microcalorimeter focal plane. We first evaluated a
reference detector--cryostat configuration. Its day-15 post-selection
background rate was $5.31\times10^{-2}\cps$ and was dominated by delayed
activation. Selected-event tracing showed that copper contributed 70.87\% of
the delayed background, principally from the 50 mK cold plate, TES-support
copper panels, and TES copper heat-sink ring. The concentration of selected
events in these near-field structures motivated evaluation of a complete
lateral-chimney configuration, in which the TES was displaced from the
cold-stage projection and surrounded by bismuth germanate (BGO) active
shielding. This configuration retained 99.75\% of the reference post-selection
effective area while reducing the day-15 background to
$1.07\times10^{-2}\cps$. Its residual background comprised \revblue{atmospheric
gamma rays} (49.8\%), delayed activation (33.3\%), and non-gamma
prompt particles (16.9\%). For a 20 d balloon observation, the top-of-atmosphere
$3\sigma$ minimum resolvable line flux was
$(2.44\pm0.22)\times10^{-5}\phcms$, a factor of 2.23 below that of the reference
configuration. The results support two design priorities for low-background
511 keV telescopes: separating the TES from activation-prone near-field
materials and their unshielded lines of sight, and increasing the solid angle
covered by active anticoincidence.

\vspace{0.7em}
\noindent\textbf{Keywords}\hspace{0.7em}Gamma-ray instrumentation, 511 keV annihilation line, transition-edge sensors, Laue lens, balloon-borne telescope, activation background

\vspace{0.3em}
\noindent\rule{\textwidth}{0.4pt}
\vspace{1.0em}

\section{Introduction}
\label{sec:introduction}

Galactic 511 keV annihilation radiation is a key observational probe of the
production, propagation, and annihilation environments of Galactic positrons.
Its spectrum contains a two-photon line near 511 keV and an ortho-positronium
continuum at lower energies. Positrons annihilate directly with electrons or
form para-positronium, producing two photons near 511 keV, whereas
ortho-positronium decays into a three-photon continuum. The line intensity,
profile, and positronium continuum therefore constrain the medium in which
positrons slow down and annihilate \cite{Prantzos2011,Jean2006}. Under an
approximate steady-state assumption, the total Galactic 511 keV flux measured
by INTEGRAL/SPI corresponds to an annihilation rate of a few
$\times10^{43}\,\mathrm{e^+\,s^{-1}}$ \cite{Prantzos2011,Kierans2020}. Proposed
positron sources include $\beta^+$-unstable nuclei synthesized in stellar
explosions, compact objects capable of producing electron--positron pairs, and
more speculative particle-physics scenarios in which annihilating or decaying
light dark matter injects low-energy positrons. Existing observations do not
uniquely determine the dominant contribution
\cite{Prantzos2011,Siegert2016AA,Siegert2016V404}.

Wide-field observations have established the diffuse morphology of the line.
INTEGRAL/SPI finds a bright central bulge and a weaker disk, and a recent
20-year analysis recovers central, broad-bulge, and disk components while the
evidence for additional correlated structures remains marginal
\cite{Knodlseder2005AllSky,Siegert2016AA,Yoneda2025}. The 2016 COSI balloon flight independently
detected the bulge and found evidence that the emission is more extended than a
single point source \cite{Kierans2020,Siegert2020COSI}. Spectroscopy also
constrains the annihilation medium and kinematics: the SPI spectrum can be
decomposed into a narrow component with about $1.3\keV$ FWHM and a broad
component with about $5.4\keV$ FWHM, and spatial variations in centroid energy
and width have been sought across the inner Galaxy
\cite{Jean2006,Siegert2019Kinematics}. The annihilation map, however, records
where positrons slow down and annihilate, not necessarily where they were
created. Propagation through a magnetized, multiphase interstellar medium means
that diffuse morphology and spectroscopy alone cannot recover the dominant
source class or the complete pre-annihilation transport history
\cite{Prantzos2011}.

Wide-field Compton and coded-mask instruments are well suited to the total
bulge and disk emission. A complementary pointed instrument is needed to test
whether the central region also contains an unresolved central enhancement,
compact sources, or measurable line-kinematic structure. A detection of a
point-like or narrowly concentrated feature would show that part of the
annihilation is spatially localized; a non-detection would constrain a compact
component without excluding source populations whose positrons escape and
annihilate diffusely. An INTEGRAL/IBIS compact-source search found no Galactic
511 keV point sources and set a $2\sigma$ upper limit of about
$1.6\times10^{-4}\phcms$ for a $\sim10$ Ms Galactic-centre exposure
\cite{DeCesare2011}. This limit sets both the comparison scale of earlier
compact-source searches and a sensitivity target for a new narrow-field
pointing instrument.

\revmark{1}Laue optics use transmission-mode Bragg diffraction to focus photons arriving
from a known direction onto a compact focal plane
\cite{Barriere2009Laue,Barriere2011LauePrototype}. When detector-related
background dominates, concentrating photons collected over a large aperture
onto a small sensitive detector area improves point-source sensitivity
\cite{Weidenspointner2006MAX}. The focal-plane detector must combine a narrow
energy response with sufficient stopping efficiency at 511 keV.
Transition-edge-sensor (TES) microcalorimeters measure the small temperature
rise produced by deposited energy through the steep resistive transition of a
superconductor and can achieve sub-keV energy resolution near 511 keV
\cite{Shirazi2023,Zhang2022TESGamma}. An array covers the focal spot, while thick
absorbers and a multilayer layout increase the full-energy-peak efficiency. The
Laue lens concentrates source photons onto a compact focal plane, and the TES
response supports a narrow line window and detailed line-profile measurements.
The resulting narrow-field focusing spectrometer trades sky coverage for a
concentrated source beam and high-resolution spectroscopy. It complements
previous wide-field measurements of the Galactic bulge and disk and supports
dedicated searches for compact sources.

Energy resolution and focusing performance alone do not determine point-source
sensitivity \cite{Shirazi2023,Hu2026}. Prompt energy deposits from the radiation
environment and delayed emission from activated detector and cryostat materials
also set the background in the science window
\cite{Gallego2025Balloon,Tian2026DIXE}.

\revmark{2}This work evaluates the high-altitude performance of a balloon-borne Laue-lens
TES-array telescope for detecting 511 keV point sources toward the Galactic
center. We classify the science-window background by physical component and
incident particle, trace it to the relevant materials, parent nuclides, and
production regions, and use this provenance to identify the dominant production
mechanisms and guide optimization of the detection system.

\subsection{Science and instrument requirements}
\label{sec:requirements}

\begingroup
\clubpenalty=10000
\widowpenalty=10000

\revmark{0}The instrument simulated here is designed to search for compact 511 keV sources
and an unresolved central-excess component in the Galactic-center region.
INTEGRAL/SPI measurements show that the 511 keV sky is dominated by diffuse
bulge and disk emission \cite{Siegert2016AA,Yoneda2025}. De Cesare's dedicated
search for compact sources with INTEGRAL/IBIS found no Galactic 511 keV point
sources and reported a $2\sigma$ upper limit for a 10 Ms exposure toward the
Galactic center \cite{DeCesare2011}. To compare that limit with the $3\sigma$
sensitivity benchmark reported here, we rescale it linearly under the Gaussian
approximation and define the reference source flux before atmospheric
attenuation as $\fref=2.4\times10^{-4}\phcms$. At each trajectory node, the
mission fold calculates atmospheric transmission from the residual column depth
at the assumed balloon altitude and the fixed $45^\circ$ instrument elevation.
The representative reference state is the day-15 trajectory node at
$38.75\,\mathrm{km}$; the adopted residual column at this node gives
$T_{\mathrm{ref}}=0.652034$. The corresponding telescope-incident reference
flux is $F'_{\mathrm{ref}}=\fref T_{\mathrm{ref}}=
1.565\times10^{-4}\phcms$. At other trajectory nodes, the incident flux is
$\fref T_{\mathrm{atm},511}(t)$, and minimum resolvable fluxes are reported
before atmospheric attenuation. Compared with the wide-field
spectroscopic and coded-mask searches conducted with INTEGRAL, the payload
modeled here uses Laue focusing for its point-source response and a
high-energy-resolution, multilayer thick-absorber TES microcalorimeter array at
the focal plane \cite{Shirazi2023,Caputo2017}.

\revdeleted{R4: the former four-requirements paragraph has been deleted in this review copy.}

The remainder of this paper is organized as follows. Section~\ref{sec:geometry}
presents the detector and cryostat mass models. Section~\ref{sec:framework}
describes the simulation framework, event selection, mission-time trajectory
fold, and response-weighted environmental screening. Section~\ref{sec:results}
presents the balloon post-selection rates, background provenance, and
sensitivity. Section~\ref{sec:environment_screening} transfers the response to
orbital and lunar-surface source fields. Section~\ref{sec:summary} summarizes
the background composition, final performance, and design direction, and
Section~\ref{sec:outlook} considers LEO-satellite and lunar-surface concepts.

\endgroup

\section{Detector and cryostat mass model}
\label{sec:geometry}

This study compares two complete detector--cryostat mass models: the reference
mass model A and mass model B, which was developed from the background analysis
of model A. Both contain the same six-layer TES core and a compact dilution
refrigerator (DR) with a mixing chamber, staged cold plates, thermal and magnetic
shields, copper supports, active scintillators, and cryogenic service structures
\cite{Shirazi2023,Gottardi2021TESReview}. They differ in the placement of the
focal plane and in the physical entrance, local tungsten structure, and active
shield arrangement. Figure~\ref{fig:mass_model} compares the two geometries with
the same camera and physical scale.

\revmark{5}The TES detector concept uses an absorber-coupled AlMn-film structure developed
in our laboratory. The transition temperature of AlMn can be tuned through its
composition and annealing process; the material is suitable for array fabrication
and has relatively low magnetic-field sensitivity
\cite{Xie2026AlMn,Vavagiakis2017AlMnMagnetic}. In the present conceptual design,
the TES thermometer is an approximately $200\,\mathrm{nm}$-thick AlMn film with
an Mn doping concentration of $2000\,\mathrm{ppm}$ on an atomic basis, whereas the $\gamma$-ray absorber is a
millimetre-scale metal body. The absorber consists of
$1.5\times1.5\times3.0\,\mathrm{mm^3}$ Ta pixels; the high atomic number of Ta
increases the interaction efficiency for 511 keV photons. At 511 keV, a single
$3.0\,\mathrm{mm}$ Ta layer has an estimated interaction efficiency of about
$49\%$, and the corresponding six-layer interaction probability is about
$98\%$. In the
current geometry, adjacent Ta array layers have a center-to-center spacing of
$12\,\mathrm{mm}$; this spacing was chosen with a subsequent veto based on
Compton-kinematic track reconstruction in mind. Ta was also selected for its low
superconducting specific heat in the cryogenic regime. \revmark{25}\revblue{Yan et al. tabulate a
volumetric heat capacity of $c_{\mathrm{vol}}=1.0\times10^{-2}\,
\mathrm{J\,m^{-3}\,K^{-1}}$ for Ta at $100\,\mathrm{mK}$
\cite{Yan2019TaHeatCapacity}. Using the Ta density adopted in the mass model,
$\rho_{\mathrm{Ta}}=16.69\times10^{3}\,\mathrm{kg\,m^{-3}}$, this corresponds to approximately
$5.99\times10^{-7}\,\mathrm{J\,kg^{-1}\,K^{-1}}$ and a heat capacity of approximately
$67.5\,\mathrm{pJ\,K^{-1}}$ for one $1.5\times1.5\times3.0\,\mathrm{mm^3}$ Ta
absorber.} This low specific heat limits
the heat-capacity penalty of a millimetre-scale absorber, allowing its stopping
efficiency to be increased while retaining the small heat capacity needed for
high calorimetric energy resolution. Ta also has a relatively
high photoelectric-to-Compton interaction ratio. From
the NIST XCOM cross sections near 511 keV, this ratio is approximately 0.731
\cite{NISTXCOM}. In the detector-response model, we use our laboratory experience
with a previously developed X-ray TES \cite{Xie2026AlMn} to adopt
$\Delta E_{\mathrm{FWHM}}\simeq420\,\mathrm{eV}$ at 511 keV as the estimated
energy resolution and detector-response post-processing input.

Because the TES film and Ta absorber differ greatly in dimensions and mass, and
because the TES film itself contributes negligibly to interactions at 511 keV,
the transport geometry uses the Ta absorber as the principal sensitive volume of
each TES pixel. The explicitly modelled sensitive volume comprises six layers,
each containing 376 Ta absorber pixels, for a total of 2256 pixels. Each pixel is
$1.5\,\mathrm{mm}$ long and wide and $3\,\mathrm{mm}$ thick. A Si substrate proxy
is placed below each Ta pixel; because detailed thermal conduction is not modelled,
Si represents the $\mathrm{SiO_2}$--$\mathrm{SiN_x}$--Si substrate. The TES film,
wiring, and readout details are not constructed layer by layer as explicit
interaction volumes, but enter the model through the assumed energy response and
service-mass proxies. Defined independently of the response FWHM, the retained science-window
analysis selection uses the 511 keV line window
$W_{511}=511\,\mathrm{keV}\pm420\,\mathrm{eV}$, or
$510.58$--$511.42\,\mathrm{keV}$.

The TES array in mass model A is mounted on copper supports immediately below the
staged cold plates. Along the optical path are successive aluminium thermal-shield
windows, a beryllium vacuum window, and a tungsten multihole collimator aligned
with the TES pixels. Boolean apertures through the surrounding shells define the
side-entry beam path, and the complete instrument frame is rotated by
$45^\circ$ so that the local aperture points $45^\circ$ above the horizon.

\revmark{6}\revmark{23}
\begingroup
\color{revisionblue}
Source tracing in mass model A identified the 50 mK mixing-chamber copper cold
plate, TES copper supports, and copper heat sink as important sources of 511-keV
science-window background. Mass model B therefore retains the six-layer TES core
but moves the focal plane laterally into a side chimney, outside the direct
projection of the staged cold plates, and surrounds it with local BGO
anticoincidence shielding. This configuration reduces both cold-end background
coupling and active-shield mass.
\par
\endgroup

\begin{figure}[t]
\centering
\includegraphics[width=0.98\linewidth]{figures/manuscript/fig01_mass_models_ab_en.pdf}
\caption{\revblue{Matched detector--cryostat views of (a) mass model A and (b)
mass model B, rendered with the same viewing angle and scale. In model A, the
six-layer TES array lies below the staged cold structures and behind a
multihole W collimator. Model B translates the same TES core into a lateral Al
side chimney, with local active BGO on the chimney side, front optical annulus,
and rear cold-port annulus. Red arrows show the nominal focused 511 keV paths;
selected outer envelopes are translucent or omitted for clarity.}}
\label{fig:mass_model}
\end{figure}

\begingroup
\color{revisionblue}
Figure~\ref{fig:mass_model} shows the principal geometric difference between the
two configurations. In mass model A, the six-layer Ta/TES focal plane lies
beneath the staged cold plates, with a 4.796-mm-thick Bi passive-shield layer in
front of the array. Mass model B translates the same focal plane laterally into
the Al side chimney, outside the direct projection of the cold plates, and
surrounds it with the chimney-side BGO shield, front optical annulus, and rear
cold-port annulus.
\par
\endgroup

\revmark{7}Both designs were simulated and analysed in the same MEGAlib/Geant4
framework, with common source, detector-response, science-window, timing, and
Compton-selection definitions and geometry-specific active anticoincidence.
The resulting performance comparison is between the two complete configurations.

\FloatBarrier

\section{Simulation framework}
\label{sec:framework}

\revmark{8}We used MEGAlib (Medium Energy Gamma-ray Astronomy Library), a
Geant4-based framework for the simulation and analysis of medium-energy
gamma-ray instruments \cite{Zoglauer2006,Agostinelli2003,Allison2016}. Its
Cosima module defines sources and detector geometry, transports particles
through the mass model, and records interactions for downstream detector-response
and event-selection analysis. Figure~\ref{fig:workflow} summarizes the end-to-end
calculation in three physical branches: focused signal, prompt atmospheric
background, and delayed activation background. The Laue optics and the
detector--cryostat system were represented by separate mass models. The
reference point source was first transported through the Laue-lens model; the
positions and directions of photons reaching the focal plane were written to a
MEGAlib EventList and used to initialize detector-side Monte Carlo transport.
For each detector--cryostat mass model, this focused-signal input and all
background sources were transported through the same complete model described
in Section~\ref{sec:geometry}.

Prompt transport and activation-production transport were run in parallel from
the same \revblue{eight-family EXPACS/PARMA atmospheric field, comprising
$\gamma$, $e^-$, $e^+$, $\mu^-$, $\mu^+$, $n$, $p$, and $\alpha$}. In
the delayed-activation branch, BUILDUP transport recorded nuclide yields,
material volumes, and production positions. These records were used to
calculate the activity accumulated during the flight and to construct the
radionuclide inventory at each reference epoch; delayed decays were then
sampled at the recorded production positions and transported through the
detector--cryostat model.

The focused signal and the \revblue{two background streams used in the rate
budget---the eight prompt atmospheric particle families and delayed
activation---} were processed with the same detector response within each
mass model. \revmark{9}The common science-window cut applied to every signal and
background stream after the detector response was
$510.58\leq E_{\mathrm{TES}}<511.42\keV$. Candidates in this window were then
subjected to the geometry-specific active-veto criterion and the common
Compton-topology test. Background events surviving all selections were
decomposed by incident particle, parent
nuclide, material, and production volume. This provenance first diagnoses mass
model A and then motivates evaluation of mass model B as a complete alternative
configuration. After both designs had been fixed, their selected signal and
component-resolved background rates were folded over the adopted trajectory to
obtain accumulated counts and unresolved-line sensitivity. The subsections
below define the source models, transport normalization, event selection, and
mission-time calculation.

\begin{landscape}
\centering
\vspace*{\fill}
\includegraphics[width=0.94\linewidth]{figures/manuscript/fig02_workflow_en_threebranch_20260830.pdf}
\captionof{figure}{\revmark{17}End-to-end simulation and analysis workflow. The focused
signal, prompt atmospheric background, and delayed activation form three
physical branches. \revblue{The prompt branch comprises eight independently
normalized particle families ($\gamma,e^-,e^+,\mu^-,\mu^+,n,p,\alpha$).} All
streams undergo the common detector response, active veto, and Compton-topology
test within each mass model. Selected-event provenance motivates the complete
model-B configuration; the fixed designs are then evaluated on the common
mission timeline.}
\label{fig:workflow}
\vspace*{\fill}
\end{landscape}

\FloatBarrier

\subsection{Laue optics and detector-coupled signal transport}
\label{sec:optics}

The simulated Laue lens uses mosaic Ge(111) crystals \cite{Barriere2009Laue,Barriere2011LauePrototype}. The optical design has a 10 m focal length, a single-ring reference response at 511\keV{}, and an effective area of $20.1\,\mathrm{cm^2}$ at 511\keV, calculated as
\begin{equation}
  A_{\mathrm{eff}}(E)=A_{\mathrm{inc}}\frac{N_{\mathrm{fp}}(E)}{N_{\mathrm{inc}}(E)}
  \label{eq:laue_aeff}
\end{equation}
where $A_{\mathrm{inc}}$ is the Monte Carlo incident sampling area, $N_{\mathrm{inc}}(E)$ is the number of incident photons, and $N_{\mathrm{fp}}(E)$ is the number of photons satisfying the focusing/focal-plane interception condition.
The mosaic-crystal implementation provides the specified angular acceptance and is used consistently throughout the optical transport.

The Laue optical branch models both focusing and photon transport through the Ge crystals.

Diffraction is implemented as a discrete application-layer process in Geant4 rather than as an efficiency applied after transport \cite{Agostinelli2003,Allison2016}. At angular offset $|\Delta\theta|$, the tabulated diffraction probability $R$ is first normalized to the unabsorbed fraction, $R_{\mathrm{eff}}=R/(1-A)$, where $A$ is the photoelectric-absorption probability. It is then converted to an equivalent mean free path $\lambda_{\mathrm{diff}}$ through
$1-\exp(-\ell/\lambda_{\mathrm{diff}})=R_{\mathrm{eff}}$, where $\ell$ is the photon path length through the crystal step; hence $\lambda_{\mathrm{diff}}=-\ell/\ln(1-R_{\mathrm{eff}})$. Geant4 samples this process together with photoelectric absorption and Compton scattering, and the shortest sampled interaction length determines which process occurs first. A photon not selected for diffraction therefore continues under the standard electromagnetic physics and may be absorbed, scattered, or transmitted.

In this work, the diffraction probability is determined from a precomputed XOP/CRYSTAL rocking curve. The angular dependences of the Ge(111) reflectivity, transmissivity, and absorption are generated with the CRYSTAL/diff\_pat module of the XOP toolkit (called through xoppylib using crystallographic data from DABAX) \cite{SanchezDelRio2011XOP} and stored as a transport lookup table during preprocessing. At each transport step inside a crystal, the angular deviation from the Bragg condition is evaluated and the corresponding diffraction probability is obtained by interpolation. The 30 arcsec FWHM mosaic spread is implemented by applying a Gaussian angular perturbation to the ideal Bragg-plane normal. When diffraction occurs, the outgoing direction is determined from the perturbed plane normal, while the photon energy is conserved. This treatment follows existing Geant4 implementations of crystal Bragg reflection \cite{Guan2023BraggReflection,Reiazi2025G4BraggReflection} and is applied to the energy range considered here.

Here the reference signal is a monochromatic 511\keV{} point source. Its sole
reference normalization is $\fref=2.4\times10^{-4}\phcms$ above the atmosphere
(Section~\ref{sec:requirements}), motivated by the dedicated INTEGRAL/IBIS
compact-source upper limit \cite{DeCesare2011} rather than by any single
observed object. It is constructed as a MEGAlib far-field source
\cite{Zoglauer2006}, with the detailed construction given in
Section~\ref{subsec:pointsource}. The signal is transported through the optics
branch and injected into the detector mass model as a parallel-beam source;
prompt cosmic-ray background and activation decays are transported directly in
the detector/cryostat mass model. The primary background budget is therefore
defined for the detector/cryostat branch, while the optics branch supplies the
transported focal-plane signal phase space.

\subsection{Source construction}
\label{sec:background}

The reference point-source EventList, the prompt atmospheric field, and the
delayed-decay source are constructed independently. \revblue{The eight prompt
particle families retain separate physical-rate normalizations}; activation production and delayed
decay remain separate transport stages. This organization preserves the
normalization and provenance needed to compare all selected background streams
under the common detector response and event selection.

\subsubsection{Prompt atmospheric source}
\label{subsec:prompt_source}

The prompt atmospheric field is derived from EXPACS/PARMA. PARMA3.0 is an
analytical model parameterized from PHITS atmospheric-shower calculations and
benchmarked against terrestrial cosmic-ray measurements, and EXPACS is its
public tabulation and software interface
\cite{Sato2015EXPACS,Sato2016PARMA4}. PHITS supplies the underlying
particle-transport physics: primary cosmic rays interact with air nuclei, the
secondary hadrons, leptons, photons, and neutrons are transported through the
atmospheric column, and the resulting fluences are parameterized by altitude,
geomagnetic cutoff, solar modulation, energy, and direction. In this work
EXPACS/PARMA provides the differential incident flux as a function of particle
species, kinetic energy, zenith angle, geographic position, altitude, cutoff
rigidity, and solar activity, and is used only to construct the external
radiation fields; detector response, vetoing, topology selection, and
mission-time counting are applied downstream. The reference source definitions
use latitude $34^\circ$, longitude $100^\circ$, altitude $38\,\mathrm{km}$,
vertical cutoff rigidity $R_c=11.6\,\mathrm{GV}$, and the 2025-08-31 solar
condition corresponding to $W=118.3$.

The prompt source includes photons, neutrons, electrons, positrons, protons,
alpha particles, and negative and positive muons. For each particle family the
EXPACS/PARMA differential spectrum is converted into a MEGAlib far-field area
source covering the full zenith range $\theta=0^\circ$--$180^\circ$ and full
azimuth range $\phi=0^\circ$--$360^\circ$. The zenith dependence is represented
by 20 differential equal-$\mu$ angular bins, where $\mu=\cos\theta$ and each
bin subtends $\Delta\Omega=0.6283185307\,\mathrm{sr}$: the first ten bins cover
down-going particles, the last ten up-going or albedo-like particles. Both
hemispheres are retained in the prompt transport and in the
activation-production transport, so upward albedo particles are not discarded
at source construction. Figure~\ref{fig:expacs_fullsphere_flux} shows the
resulting full-sphere flux field.

\begin{figure}[t]
\centering
\includegraphics[width=\linewidth]{figures/manuscript/fig03_corrected_kev_sources.pdf}
\caption{EXPACS/PARMA atmospheric particle field used for prompt transport and
activation production. (a) Fluxes of the eight particle families integrated over
energy and solid angle in the down-going and up-going (albedo) hemispheres.
(b) Integrated fluxes in the 20 equal-$\mu$ zenith bins ($\mu=\cos\theta$) for
each particle family. The environment is latitude $34^\circ$, longitude
$100^\circ$, altitude $38\,\mathrm{km}$, $R_c=11.6\,\mathrm{GV}$, and $W=118.3$.}
\label{fig:expacs_fullsphere_flux}
\end{figure}

The full-sphere integrated fluxes in the adopted atmospheric tables are 4.79966, 0.462239, 0.199002,
0.116981, 0.112301, 0.0114871, 0.00496893, and
$0.00557001\,\mathrm{cm^{-2}\,s^{-1}}$ for $\gamma$, $n$, $e^-$, $e^+$, $p$,
$\alpha$, $\mu^-$, and $\mu^+$, respectively. All families use 20 equal-$\mu$
angular bins and a 60 cm far-field start sphere.

\revdeleted{R10: the former source-energy unit-conversion paragraph has been deleted in this review copy.}

If $I_j(E,\mu)$ denotes the EXPACS/PARMA differential intensity for family $j$,
the flux carried by equal-$\mu$ bin $m$ is
\begin{equation}
 \Phi_{j,m}=2\pi\int_{\mu_m}^{\mu_{m+1}}\!\mathrm d\mu
             \int I_j(E,\mu)\,\mathrm dE,
 \qquad \Delta\mu=0.1,
 \label{eq:source_bin_flux}
\end{equation}
so every bin subtends $\Delta\Omega=2\pi\Delta\mu=0.62831853\,\mathrm{sr}$.
The first ten bins cover the down-going hemisphere and the second ten cover the
up-going or albedo hemisphere; azimuth is complete in both cases.

All eight particle families use a far-field area source of radius
$R=60\,\mathrm{cm}$. For each sampled direction, starting positions are drawn
uniformly from the projected disk perpendicular to that direction, and momenta
point into the mass model; the geometric rate factor is the projected area
$\pi R^2$. Together, the 20 bins cover $4\pi$ incident solid angle. The same
atmospheric-field definition is used separately for prompt transport and
activation-production calculations.

\begingroup
\color{revisionblue}
\noindent\textbf{[R14]} Each accepted production job for particle family $j$
contains the complete 20-bin angular mixture; the angular bins are not treated
as independent exposures. The family source rate and the merged equivalent
exposure are therefore
\begin{equation}
 \dot N^{\mathrm{src}}_{j}=\pi R^2\sum_{m=1}^{20}\Phi_{j,m},
 \qquad
 T^{\mathrm{eq}}_{g,j}=\sum_{q\in\mathcal Q_{g,j}}\mathrm{TT}_{g,j,q}
 \simeq
 \frac{\sum_q N^{\mathrm{sim}}_{g,j,q}}{\dot N^{\mathrm{src}}_{j}},
 \qquad
 w^{\mathrm{pr}}_{g,j}=\frac{1}{T^{\mathrm{eq}}_{g,j}},
 \label{eq:prompt_equivalent_exposure}
\end{equation}
where $q$ indexes independent mergeable jobs and $\mathrm{TT}_{g,j,q}$ is the
validated physical runtime recorded in the corresponding receipt. Increasing
the number of transported primaries extends $T^{\mathrm{eq}}$ and reduces Monte
Carlo uncertainty without changing the physical source intensity.

\revdeleted{R24: The catalogue-rate and common-time Poisson-sampling equation and its expanded explanation were deleted here.}
\endgroup

\subsubsection{Delayed activation source}
\label{subsec:delayed_source}

The delayed source is built from a separate activation-production transport. It
is not counted as immediate detector background; instead, it records the isotope,
material volume, incident family, and production position of each radioactive
product. Ground-state half-lives are checked against NUBASE2020
\cite{Kondev2021NUBASE}. For isotope $k$,
\begin{equation}
  A_k(t_{\mathrm{flight}})=P_k\left[1-\exp(-\lambda_k t_{\mathrm{flight}})\right],
  \qquad \lambda_k=\frac{\ln2}{t_{1/2,k}},
\end{equation}
where $P_k$ is the production rate and the reference state is
$t_{\mathrm{flight}}=15\,\mathrm{d}$.

\begingroup
\color{revisionblue}
\noindent\textbf{[R11]} In the standard output configuration used here,
MEGAlib/Cosima's compressed \texttt{sim.gz} event files do not retain the
complete radioactive-product production information required to reconstruct
the spatial distribution of the delayed source. Because Cosima uses Geant4 as
its transport engine, we added a recording hook to \texttt{SteppingAction} to
store the production coordinate, material volume, nuclide, and excitation
state when each radioactive product was created. These fields were then used
to match the day-15 inventory to the production records and to organize the
delayed source by incident family, material category, production volume, and
coordinate.

Retaining the production positions is essential because the ability of a
delayed particle to escape the surrounding material, trigger the
geometry-specific active anticoincidence, or pass the Compton-topology
selection depends on where it was produced. Decay positions are sampled from
the matched production records in proportion to activity, preserving both the
spatial distribution and total activity. The activation-source positions in
the two mass models, together with the parent nuclides contributing to the TES
511 keV response window, are shown in
Fig.~\ref{fig:activation_source_positions}.
\endgroup

\revdeleted{R12: the redundant method-level statement of finite post-selection event counts has been deleted in this review copy.}

\begingroup
\color{revisionblue}
\noindent\textbf{[R15]} For each incident-particle family, each mass model
transports $N^{\mathrm{dec}}_{g,j}=10^6$ delayed decays. At the day-15 reference
epoch, the constant-activity equivalent exposure and per-event rate weight of
that family are
\begin{equation}
 T^{\mathrm{eq,del}}_{g,j}(t_{\mathrm{ref}})
 =\frac{N^{\mathrm{dec}}_{g,j}}{A_{g,j}(t_{\mathrm{ref}})},
 \qquad
 w^{\mathrm{del}}_{g,j}(t_{\mathrm{ref}})
 =\frac{A_{g,j}(t_{\mathrm{ref}})}{N^{\mathrm{dec}}_{g,j}}
 =\frac{1}{T^{\mathrm{eq,del}}_{g,j}(t_{\mathrm{ref}})},
 \quad t_{\mathrm{ref}}=15\,\mathrm{d}.
 \label{eq:delayed_equivalent_exposure}
\end{equation}
\revmark{24}Here $T^{\mathrm{eq,del}}$ is the Monte Carlo equivalent sampling time for the
day-15 activity mixture, rather than the elapsed flight time or the BUILDUP
irradiation time; the mission fold updates the day-15 reference weights at each
time node using the parent-nuclide activity curves.
\endgroup

\begingroup
\color{revisionblue}
\noindent\textbf{[R13]} For activation background near 511 keV, the final
events in both models are primarily associated with copper structures in the
cryogenic assembly. The source decomposition for mass model A identifies
copper near the TES cold stages and mixing chamber as an important
contribution. Guided by this diagnosis, mass model B places the focal plane in
a lateral chimney to reduce its geometric coupling to the principal copper
activation regions and rearranges local BGO along the cold-stage direction,
thereby increasing the probability of a coincident anticoincidence tag. Under
the common source, detector-response, and selection framework, the complete
mass-model-B configuration yields a lower delayed-background rate than mass
model A; Section~\ref{sec:results} examines the source--geometry relationship
further.
\endgroup

\begingroup
\color{revisionblue}
\begin{table}[!htbp]
\color{revisionblue}
\centering
\captionsetup{font={color=revisionblue}}
\caption{\textbf{[R16]} Particle-family transport statistics and rate
normalization for the two mass models. $N_{\mathrm{tr}}$ and $N_{\mathrm{dec}}$
are transported prompt primaries and delayed decays, respectively;
$T_{\mathrm{eq}}$ is the summed physical exposure of accepted independent
production batches; and $N_{\mathrm{det+}}$ counts templates with a positive
raw deposit in the TES or an active shield. $R_{\mathrm{cat}}$ is the
detector-positive occupancy rate actually supplied to common-time Poisson
sampling. It is evaluated before detector-response smearing, the pixel
threshold, science window, active veto, and Compton selection. In panel (b),
$T_{\mathrm{eq,del}}=N_{\mathrm{dec}}/A_{g,j}(15\,\mathrm d)$, as defined in
Equation~(\ref{eq:delayed_equivalent_exposure}).}
\label{tab:background_source_model}
\scriptsize
\setlength{\tabcolsep}{5.0pt}
\renewcommand{\arraystretch}{0.92}
\noindent\textit{(a) Prompt atmospheric-particle transport}\par\vspace{0.20em}
\begin{tabular}{clrrrr}
\toprule
Model & Family & $N_{\mathrm{tr}}$ & $T_{\mathrm{eq}}$ (s)
& $N_{\mathrm{det+}}$ & $R_{\mathrm{cat}}$ ($\mathrm{s^{-1}}$) \\
\midrule
A & $\gamma$ & $20{,}315{,}706$ & 374.249513 & $6{,}372{,}612$ & $1.635799\times10^{4}$ \\
A & $n$ & $1{,}164{,}960$ & 223.112900 & $403{,}599$ & $1.808945\times10^{3}$ \\
A & $e^-$ & $675{,}057$ & 299.757900 & $287{,}509$ & 959.137357 \\
A & $e^+$ & $394{,}119$ & 298.875770 & $182{,}318$ & 610.012648 \\
A & $p$ & $380{,}829$ & 299.752500 & $172{,}734$ & 576.255411 \\
A & $\alpha$ & $39{,}050$ & 298.957500 & $18{,}657$ & 62.406864 \\
A & $\mu^-$ & $16{,}607$ & 298.151700 & $7{,}529$ & 25.252246 \\
A & $\mu^+$ & $19{,}371$ & 308.674500 & $8{,}629$ & 27.955014 \\
\addlinespace[0.25em]
B & $\gamma$ & $112{,}421{,}262$ & 2071.102900 & $4{,}189{,}660$ & $1.939983\times10^{3}$ \\
B & $n$ & $8{,}363{,}786$ & 1599.414700 & $325{,}382$ & 203.438170 \\
B & $e^-$ & $3{,}602{,}914$ & 1598.881500 & $136{,}738$ & 85.521035 \\
B & $e^+$ & $2{,}119{,}548$ & 1600.953780 & $139{,}963$ & 87.424760 \\
B & $p$ & $2{,}036{,}710$ & 1602.777690 & $123{,}624$ & 77.131096 \\
B & $\alpha$ & $209{,}074$ & 1607.562620 & $18{,}285$ & 11.374363 \\
B & $\mu^-$ & $91{,}020$ & 1608.091082 & $4{,}605$ & 2.863644 \\
B & $\mu^+$ & $102{,}306$ & 1612.825350 & $5{,}268$ & 3.266318 \\
\bottomrule
\end{tabular}

\vspace{0.45em}
\noindent\textit{(b) Delayed-decay transport at day 15}\par\vspace{0.20em}
\begin{tabular}{clrrrr}
\toprule
Model & Family & $N_{\mathrm{dec}}$ & $T_{\mathrm{eq,del}}$ (s)
& $N_{\mathrm{det+}}$ & $R_{\mathrm{cat},15}$ ($\mathrm{s^{-1}}$) \\
\midrule
A & $\gamma$ & $1.0\times10^6$ & $1.516919\times10^5$ & $884{,}678$ & 5.832072 \\
A & $n$ & $1.0\times10^6$ & 3030.110 & $893{,}664$ & 294.927907 \\
A & $e^-$ & $1.0\times10^6$ & $1.497281\times10^6$ & $853{,}592$ & 0.570095 \\
A & $e^+$ & $1.0\times10^6$ & $3.385589\times10^5$ & $904{,}775$ & 2.672430 \\
A & $p$ & $1.0\times10^6$ & 1380.694 & $916{,}824$ & 664.031386 \\
A & $\alpha$ & $1.0\times10^6$ & 3477.763 & $918{,}773$ & 264.185032 \\
A & $\mu^-$ & $1.0\times10^6$ & $4.034086\times10^6$ & $871{,}246$ & 0.215971 \\
A & $\mu^+$ & $1.0\times10^6$ & $7.933795\times10^{14}$ & $1{,}000{,}000$ & $1.260431\times10^{-9}$ \\
\addlinespace[0.25em]
B & $\gamma$ & $1.0\times10^6$ & $4.921159\times10^5$ & $396{,}106$ & 0.804904 \\
B & $n$ & $1.0\times10^6$ & $1.459182\times10^4$ & $554{,}831$ & 38.023436 \\
B & $e^-$ & $1.0\times10^6$ & $3.754719\times10^6$ & $433{,}022$ & 0.115327 \\
B & $e^+$ & $1.0\times10^6$ & $1.480682\times10^6$ & $531{,}844$ & 0.359189 \\
B & $p$ & $1.0\times10^6$ & 9476.321 & $694{,}056$ & 73.241080 \\
B & $\alpha$ & $1.0\times10^6$ & $2.916849\times10^4$ & $711{,}199$ & 24.382444 \\
B & $\mu^-$ & $1.0\times10^6$ & $1.461788\times10^7$ & $436{,}535$ & 0.029863 \\
B & $\mu^+$ & $1.0\times10^6$ & $2.521327\times10^8$ & $1{,}867$ & $7.404830\times10^{-6}$ \\
\bottomrule
\end{tabular}

\vspace{0.35em}
\parbox{0.97\linewidth}{\raggedright For the \revblue{atmospheric-$\gamma$ prompt
catalogue}, $R_{\mathrm{cat}}=\sum_e w_e$ includes per-event spectral-importance
weights; the seven non-$\gamma$ prompt families use the constant
$1/T_{\mathrm{eq}}$ weight. Delayed weights are constant within a family at day
15, while the mission fold evolves individual parent nuclides rather than
rescaling an entire family by one factor.}
\end{table}
\endgroup

\begin{figure}[p]
\centering
\includegraphics[width=0.96\linewidth]{figures/manuscript/fig_activation_positions_model_a_en.pdf}
\vspace{0.35em}

\includegraphics[width=0.96\linewidth]{figures/manuscript/fig_activation_positions_model_b_en.pdf}
\caption{Delayed activation-production positions in mass model A (top) and mass
model B (bottom). In each row, the left panel shows the full mass model and the
right panel enlarges the TES region. Blue points denote all transported delayed
activation-source positions; coloured points denote activation parent nuclides
that produce background in the TES $510.58$--$511.42$ keV response window.}
\label{fig:activation_source_positions}
\end{figure}

\FloatBarrier

\subsubsection{Reference point-source construction}
\label{subsec:pointsource}

\begingroup
\color{revisionblue}
\revmark{18}The focused-signal branch uses an on-axis monochromatic 511\keV{}
far-field point source, represented by a parallel beam at the lens entrance.
The flux convention and atmospheric attenuation of the reference source are
defined in Section~\ref{sec:requirements}; this subsection specifies only the
focal-plane photon phase space and the detector effective-area calculation.

In the optical transport, incident photons are sampled uniformly over the lens
aperture $A_{\mathrm{inc}}$ defined by Equation~\eqref{eq:laue_aeff}. The
focal-plane response is set by the 30 arcsec crystal mosaic spread, the
XOP/CRYSTAL rocking curve, and the 10 m focusing geometry. Three independent
optical samples contain $1.5\times10^5$ incident photons and produce 37,256
focal-plane crossings. Of these, $N_{\mathrm{fp}}=37{,}194$ fall within the
$r=1.898\,\mathrm{cm}$ beryllium entrance window. The corresponding encircled
spot radii are $r_{90}=1.03\,\mathrm{cm}$ and $r_{99}=1.25\,\mathrm{cm}$. The
position and direction of each accepted photon are written to a common
focal-plane EventList and transported separately through mass models A and B.

This EventList represents the optical effective area
$A_{\mathrm{fp}}\equiv A_{\mathrm{eff}}(511\keV)=20.08476\,\mathrm{cm^2}$.
Each equally weighted focal-plane photon therefore carries the area weight
\begin{equation}
 a_{\mathrm{ray}}=\frac{A_{\mathrm{fp}}}{N_{\mathrm{fp}}}.
 \label{eq:focal_eventlist_area_weight}
\end{equation}
For mass model $g$, let $N_{\mathrm{sel},g}$ be the number of focal-plane
photons retained after detector transport, response convolution, science-window
selection, the geometry-specific isolated-event active-anticoincidence
criterion, and the Compton topology/aperture selection. Both models use the
same Compton reconstruction algorithm, while aperture compatibility is tested
against the physical entrance of the corresponding model. The final on-axis
isolated-event effective area is
\begin{equation}
 A_{\mathrm{eff},g}
 =\sum_{i\in\mathrm{sel},g}a_{\mathrm{ray}}
 =A_{\mathrm{fp}}\frac{N_{\mathrm{sel},g}}{N_{\mathrm{fp}}}.
 \label{eq:selected_signal_effective_area}
\end{equation}
\revdeleted{R24: The qualifying paragraphs on accidental overlap, off-axis response, and attitude jitter were deleted here.}
\endgroup

\revdeleted{R18: Repeated reference-flux, atmospheric-transmission, and day-15
performance statements were removed from this Methods subsection. Flux folding
remains in the mission-time calculation, and the A/B values remain in Results.}

\subsection{Detector response, event selection, and veto definitions}
\label{sec:selection}

The common-time calculation maps the prompt atmospheric and delayed-activation
background catalogues onto one explicit time axis before the active-veto and
Compton selections are applied. \revblue{The eight prompt particle families
retain independent physical-rate normalizations}. Focused-signal events are inserted as
conditional probes of accidental-coincidence loss rather than added to the
background budget.

\subsubsection{Time axis and rate normalization}
\label{subsec:time_axis}

In a stationary radiation environment, mutually independent particle events
with a constant mean arrival rate have Poisson-distributed counts in a fixed
time interval. Conditional on the event count, their arrival times are
uniformly distributed within that interval. This provides a natural temporal
representation for rate-normalized Monte Carlo catalogues: the catalogue
supplies the energy deposits, positions, and topologies, while the Poisson
process determines how often and when those events occur on the physical time
axis.

The radiation environment evolves along the balloon trajectory, but much more
slowly than the $1\,\mu\mathrm{s}$ coincidence time scale adopted in this study.
Within each trajectory node, the event streams can therefore be approximated as
independent Poisson processes with constant arrival rates. The focused signal,
prompt atmospheric background, and delayed-activation background are generated
by separate transport calculations and do not share a physical time axis,
whereas the geometry-specific active anticoincidence and Compton-topology selections depend on all
energy deposits within the same time window. Evaluating accidental coincidences
therefore requires these catalogues to be mapped onto a common time axis.

The three transported streams acquire per-record rate weights $r_{km}$ during
the source-normalization stage of Section~\ref{sec:background}, where
$k\in\{\mathrm{signal},\mathrm{prompt},\mathrm{delayed}\}$ labels the catalogue
and $m$ labels a transported record. The total rate of catalogue $k$ is
$R_k=\sum_m r_{km}$. For a reference interval of duration $T$, the number of
instances drawn from catalogue $k$ is
\begin{equation}
 N_k\sim\mathrm{Pois}(R_kT).
 \label{eq:poisson_superposition}
\end{equation}
Records are sampled with replacement with probabilities $r_{km}/R_k$ and are
assigned independent uniform times in $[0,T]$. The resulting sequence preserves
both the physical rate and the energy, direction, and detector-response
distribution of the event stream.

The three draws are merged and ordered by arrival time. Events separated by no
more than $\tau=1\,\mu\mathrm{s}$ are linked transitively into one coincidence
candidate group $c$. All energy deposits in the group are then subjected jointly
to the science-window, geometry-specific active-anticoincidence, and Compton-topology selections.
Figure~\ref{fig:poisson_time_axis_schematic} summarizes this construction.

\begin{figure}[H]
\centering
\includegraphics[width=\linewidth]{figures/manuscript/fig04a_poisson_time_axis_schematic_threebranch_20260830.pdf}
\caption{Schematic construction of the common-time coincidence axis. (a) Each
transported stream carries per-record rate weights and a total rate.
(b) Independent Poisson draws are merged and time-ordered on a common axis.
(c) Events separated by no more than $\tau=1\,\mu\mathrm{s}$ form one
coincidence candidate group. Colours identify catalogue membership; event
positions and densities are schematic. Prompt and delayed streams define the
background axis; focused signal is inserted as a conditional probe.}
\label{fig:poisson_time_axis_schematic}
\end{figure}

This calculation is performed at five reference times on days 0, 5, 10, 15,
and 20. At each reference time, the physical rates of the \revblue{eight prompt
particle families and delayed activation} set the
Poisson draws used to construct a common background time axis for that
reference environment. Events separated by no more than $1\,\mu\mathrm{s}$ form
one coincidence group and jointly undergo the science-window, geometry-specific
active-anticoincidence, and Compton-topology selections. Focused-signal events are
also inserted individually into the corresponding background axis to quantify
signal loss from accidental coincidences.

Instrument performance is then evaluated along the long-duration balloon
trajectory, with the 20 d flight divided into 6 h intervals. Particle fluxes,
atmospheric transmission of the signal, and the activation inventory evolve with flight
time and position, thereby changing both the event density on the common time
axis and the accidental-coincidence probability. The background corrections
and signal-survival fractions obtained at the five reference times are applied
along the full trajectory. Section~\ref{sec:mission_time_fold} specifies the
trajectory model, temporal interpolation, and cumulative-count calculation.

\subsubsection{Detector response, science window, and weighted rates}
\label{subsec:detector_response}

Before the science-window and anticoincidence decisions, TES energy deposits in
each candidate group are combined by pixel. Multiple deposits in the same pixel
are summed to give its true deposited energy, and their energy-weighted position
represents the interaction location. The expected single-pixel energy resolution
is $\Delta E_{\mathrm{FWHM}}=420\,\mathrm{eV}$; TES measurement fluctuations are
approximated by a Gaussian response with this FWHM. Pixel hits with measured
energies below $0.300\keV$ are removed. Mass models A and B use the same response
and threshold settings.

The measured energies of all retained TES pixel hits are summed to give the total
measured TES energy of the event, which is used for the 511\,keV science-window
selection. The energy, position, layer, and multiplicity of each pixel hit remain
available for the subsequent Compton-topology selection. BGO deposits are
accumulated separately, are not included in the TES total, and enter the active
anticoincidence decision in the next subsection.

The primary analysis uses the common science-window condition
$510.58\leq E_{\mathrm{TES}}<511.42\keV$. This
high-containment reference window was fixed in advance from the expected energy
resolution. It is applied identically to the focused signal and every background
component and is held fixed between the two mass models, rather than being tuned
separately for each geometry.

\revdeleted{R24: The qualifying statement on intrinsic line width and re-optimization of the science window was deleted here.}

The source catalogues have different physical fluxes, transport exposures, and
source-level reweighting. Each event therefore carries a physical-rate weight,
and the signal or background rate at a given selection stage is obtained by
summing the weights of events that pass that stage. Raw record counts quantify
Monte Carlo transport support rather than physical rate. Finite-transport
statistical uncertainties are propagated from the actual weight distribution,
while focused-signal acceptance is treated binomially using the common focal-plane
sample. Long-duration rate propagation along the mission trajectory is described
in Section~\ref{sec:mission_time_fold}; the response-convolved spectra through
the active-veto and topology/aperture stages are reported with the cut-flow results.

\subsubsection{Active anticoincidence gate}

Active anticoincidence rejects events that produce a 511\,keV candidate in the
TES while also depositing energy in the surrounding active shield. Passive
shielding attenuates incident radiation, whereas active BGO provides
time-coincident information for event-by-event discrimination. A focused
511\,keV photon reaches the TES mainly through the optical channel. Atmospheric
charged particles, secondaries from hadronic interactions, and some
\revblue{high-energy atmospheric-$\gamma$ cascades} are more likely to deposit energy in both the TES
and the surrounding active volumes. Coincident BGO energy can therefore
identify background candidates in the science window.

The active gate uses the common time axis and $1\,\mu\mathrm{s}$ coincidence
groups defined in Section~\ref{subsec:time_axis}. Within each candidate group,
TES deposits are first combined by pixel and then passed through the response
defined in Section~\ref{subsec:detector_response}; the science-window decision
uses the sum of the measured energies of all retained pixels. Simulated energy
deposits in all BGO volumes in the same group are summed as
$E_{\mathrm{BGO}}$. This active-volume sum is used only for anticoincidence and
does not contribute to the TES event energy.

The active system of mass model A comprises a BGO side shell, bottom cap, and
top annulus. The active system of mass model B comprises the chimney side
shield, front optical annulus, and rear cold-port annulus. Both mass models use
the same group-level passing condition,
\[
 E_{\mathrm{BGO}}<50\keV .
\]
The threshold is applied to the summed BGO deposit over the entire coincidence
group; a sum equal to or above $50\keV$ vetoes the group.
Within a given mass model, the same group-level criteria are applied to the
focused signal, \revblue{eight-family prompt background and delayed
activation}.

The groupwise sums include both correlated TES--shield deposits from one
incident particle and accidental overlap of otherwise independent hits whose
arrival times satisfy the $1\,\mu\mathrm{s}$ linking rule. A focused photon can
therefore be rejected when an unrelated active-shield hit enters the same
coincidence group. The associated loss is described by the accidental
signal-survival fraction defined in Section~\ref{subsec:time_axis} and included
in the mission-time fold. Candidates that pass the active gate proceed to the
Compton-topology and entrance-aperture consistency selection.

The analysis uses simulated energy deposits recorded by Cosima and applies
$50\keV$ as the common offline anticoincidence threshold.

\subsubsection{Compton topology and entrance-aperture consistency}

Active anticoincidence removes candidates with energy deposits above the adopted
threshold in the surrounding scintillators, but it cannot distinguish a focused photon from a
background event whose interactions remain within the TES array. Either can
leave a total energy near 511\,keV in the focal plane. A focused photon must,
however, arrive through the side-entry optical channel. For an event with two
or more separated TES hits, the division of energy among the hits and their
relative positions constrain the possible arrival directions. We therefore test
whether each event is compatible with the entrance aperture. The test uses the
kinematic cone construction of a Compton camera
\cite{Zoglauer2006}, but only as an aperture-consistency veto, not for imaging
or source localization.

The six-layer pixelated TES focal plane provides the information needed for
this test at pixel scale. For the retained hits defined in
Section~\ref{subsec:detector_response}, the response model supplies a measured
energy, while the event catalogue preserves the TES layer and pixel position.
Hits in different pixels or layers therefore provide both an energy division
and a finite spatial lever arm. The catalogue does not determine the interaction
order or the position within a pixel. A single-hit event consequently carries
no Compton-direction information and is retained as a calorimetric line
candidate; multi-hit events proceed to the consistency test below.

For an assumed interaction order, let a candidate of total energy $E_0$ deposit
$E_1$ at the first representative position $\mathbf r_1$ and then propagate
with $E_{\mathrm{sc}}=E_0-E_1$ to a second position $\mathbf r_2$. The measured
energies give the Compton scattering angle
\begin{equation}
  \cos\theta_C =
  1 - m_ec^2\left(\frac{1}{E_{\mathrm{sc}}}
  - \frac{1}{E_1+E_{\mathrm{sc}}}\right).
\end{equation}
The vector $\mathbf r_2-\mathbf r_1$ fixes the scattered-photon direction
$\hat{\mathbf u}_{12}$, but not the incident direction. Together,
$\hat{\mathbf u}_{12}$ and $\theta_C$ define a cone of allowed incident
directions with apex $\mathbf r_1$ and axis $-\hat{\mathbf u}_{12}$. The event
is aperture compatible if this cone intersects a circular analysis disk in the
side-entry plane. The disk radii are $1.898\,\mathrm{cm}$ for mass model A and
$1.900\,\mathrm{cm}$ for mass model B. These disks define the geometric region
used by the selection; they are not the complete physical windows or collimator
openings.

The remaining position uncertainty is handled explicitly. The catalogue gives
an analysis position for each hit, but does not resolve the interaction point
within the pixel. We represent that uncertainty by an axis-aligned box centred
on the recorded position, with half-lengths $0.075$, $0.075$, and
$0.150\,\mathrm{cm}$ along the analysis $x$, $y$, and $z$ axes. Its eight
vertices and centre give nine representative positions, not nine additional
measurements. For an assumed order of the first two hits, the two boxes produce
$9\times9=81$ position pairs and hence 81 candidate cones. Each cone is sampled
at 24 uniformly spaced azimuths and extended forward to the entrance plane. The
cone intersects the disk if a sampled point lies inside it or if a segment
joining adjacent samples crosses it. The test is deliberately permissive: one
intersection among the 81 candidate cones is enough to retain that order.
Figure~\ref{fig:compton_aperture_consistency} summarizes this construction.

\begin{figure}[H]
\centering
\includegraphics[width=\linewidth]{figures/fig_compton_aperture_consistency_en.pdf}
\caption{Schematic of the Compton-cone/aperture-consistency veto. (a) The first
two TES hits define a reconstruction cone from their measured energies and
representative positions; the cone is a set of allowed source directions, not a
unique photon trajectory. (b) Each hit is represented by the eight vertices and
centre of an axis-aligned analysis box, giving $9\times9=81$ position hypotheses
for a specified order. (c) The cone trace in the entrance plane is approximated
by 24 azimuthal samples and their connecting segments. An intersection with the
analysis-aperture disk retains the event, whereas a valid reconstruction that
misses the disk is vetoed. Events without a valid cone remain in the rate
accounting as unreconstructed. The schematic is not to scale.}
\label{fig:compton_aperture_consistency}
\end{figure}

The interaction order is treated separately from the within-pixel position
uncertainty. For a two-hit event, both orders are tested and either compatible
order is sufficient. For an event with at least three hits, all orders are first
enumerated using the recorded analysis positions. For a candidate sequence of
$n$ hits $\{(\mathbf r_j,\epsilon_j)\}_{j=1}^{n}$, the Compton sequence residual
is
\begin{equation}
 Q=\sum_{i=2}^{n-1}\left(\cos\theta_{C,i}-
 \hat{\mathbf u}_{i-1,i}\!\cdot\!\hat{\mathbf u}_{i,i+1}\right)^2,
\end{equation}
where $\theta_{C,i}$ is calculated from the energy $\epsilon_i$ deposited at
hit $i$ and the remaining energy $\sum_{k=i+1}^{n}\epsilon_k$; the two unit
vectors describe the adjacent flight segments. Orders without a physical
Compton solution do not enter the ranking. We select the order with the smallest
$Q$; if two orders are tied, we choose the one with the longer first lever arm.
Only the first-scatter cone of this selected order receives the 81-position
aperture test.

The final decision is conservative. A reconstructable multi-hit event is vetoed
only when the cone or cones tested under the applicable ordering rule miss the
analysis aperture. If no candidate order has a physically valid solution, the
event is marked as unreconstructed and retained in the final rate rather than
rejected. In this way, failure to reconstruct an event is not itself treated as
evidence that the event is background.

\revmark{21}\revmark{24}\revblue{The veto efficiency obtained with the present Compton-trajectory analysis should be regarded as a conservative lower bound. Offline analysis may achieve higher veto efficiency by optimizing the trade-off between signal retention and background rejection. Increasing the spacing between adjacent TES layers would lengthen the Compton-reconstruction lever arm of multilayer events and could further improve veto efficiency.}

\subsection{From the day-15 reference result to mission-time evolution}
\label{sec:mission_time_fold}

The detector transport provides component responses for the day-15 reference
environment, but the residual atmospheric depth, geomagnetic environment, and
activation inventory evolve during a multi-day float
\cite{Gallego2025Balloon,Kierans2017}. The mission projection therefore does not
extend the total day-15 background as a constant rate. Instead, it holds the
mass-model response fixed and evolves the focused signal, the eight prompt
particle families, and delayed activation separately.

The synthetic reference trajectory spans 0--20 d in 81 nodes separated by
0.25 d. Its altitude ranges from 36.25 to 41.25 km, its latitude from
$33.75^\circ$ to $34.25^\circ$ N, and its longitude from $99.75^\circ$ to
$100.25^\circ$ E. At day 15, the reference state is 38.75 km,
$34^\circ$ N, and $100^\circ$ E. The instrument axis remains fixed at
$45^\circ$ elevation. Altitude changes affect both the incident
particle field and the atmospheric transmission at 511\,keV; latitude and
longitude provide a controlled geomagnetic modulation. The trajectory is used
to compare mass models A and B under the same assumptions. It is not balloon
telemetry and does not include source visibility, Earth occultation, or
repointing history.

For a specified top-of-atmosphere 511\,keV line flux, the focused-signal rate
away from day 15 is scaled by the atmospheric transmission,
\begin{equation}
 S^{\mathrm{dir}}_g(t)
 =S^{\mathrm{dir}}_g(t_{15})
 \frac{T_{\mathrm{atm},511}(t)}{T_{\mathrm{atm},511}(t_{15})}
 =F_{\mathrm{TOA},511}A_{\mathrm{eff},g}T_{\mathrm{atm},511}(t),
 \label{eq:mission_signal}
\end{equation}
where $g$ denotes the mass model and $A_{\mathrm{eff},g}$ includes the detector
response and the single-event acceptance of the active anticoincidence and
Compton topology/aperture selections. In the plane-parallel approximation,
$T_{\mathrm{atm},511}(t)=\exp[-\mu_{\mathrm{air}}X(t)/\sin\epsilon]$. The
day-15 vertical residual column is $3.461\,\mathrm{g\,cm^{-2}}$; with
$\epsilon=45^\circ$ and an effective dry-air mass attenuation coefficient of
$0.08736\,\mathrm{cm^2\,g^{-1}}$ \cite{NISTXCOM}, this gives
$T_{\mathrm{atm},511}(t_{15})=0.652034$. The fold recalculates this transmission
at every trajectory node.

Prompt background is scaled separately for each incident family $a$,
\begin{equation}
 R_{\mathrm{prompt},g,a}(t)
 =R_{\mathrm{prompt},g,a}(t_{15})
 \frac{\Phi_a(t)}{\Phi_a(t_{15})}
 =\Phi_a(t)K_{g,a},
 \label{eq:prompt_source_response}
\end{equation}
where $\Phi_a(t)$ is the environmental flux at the trajectory node and
$K_{g,a}$ is the selected response coefficient from the day-15 reference
transport. The seven non-$\gamma$ prompt families and the \revblue{atmospheric
$\gamma$ field} use separate flux scales. Only the component normalization
changes: the within-family energy and angular templates and the detector
response remain those of the reference transport. The fold therefore does not
repeat full particle transport at the 81 nodes.

Delayed activation cannot be obtained by multiplying the total day-15 delayed
rate by one scale because the retained nuclides have different production
histories and half-lives. The production rate of source-parent $k$ from family
$a$ is
\begin{equation}
 P_{g,a,k}(t)
 =P_{g,a,k}(t_{15})\frac{\Phi_a(t)}{\Phi_a(t_{15})}
 =\Phi_a(t)Y_{g,a,k},
 \label{eq:activation_production}
\end{equation}
where $Y_{g,a,k}$ is fixed by the mass-model-specific production catalogue.
The production rate is interpolated linearly between adjacent nodes and
convolved exactly with radioactive decay,
\begin{equation}
 \begin{aligned}
 N_{g,a,k}(t_{\ell+1})={}&N_{g,a,k}(t_\ell)e^{-\lambda_k\Delta t_\ell}\\
 &+\int_0^{\Delta t_\ell}P_{g,a,k}(t_\ell+u)
 e^{-\lambda_k(\Delta t_\ell-u)}\,\mathrm du,\\
 A_{g,a,k}(t_\ell)={}&\lambda_kN_{g,a,k}(t_\ell).
 \end{aligned}
 \label{eq:delayed_inventory}
\end{equation}
The selected delayed rate is
\begin{equation}
 R_{\mathrm{delayed},g}(t)=
 \sum_a\sum_k A_{g,a,k}(t)\varepsilon_{g,a,k}.
 \label{eq:delayed_selected}
\end{equation}
Advancing each family--parent component preserves its half-life and selected
decay response. The initial condition is $N_{g,a,k}(0)=0$, so the calculation
does not include ground pre-activation or inventory accumulated during ascent
before the first trajectory node.

The total direct background before cross-event grouping is
\begin{equation}
 B^{\mathrm{dir}}_g(t)
 =\sum_aR_{\mathrm{prompt},g,a}(t)+R_{\mathrm{delayed},g}(t).
 \label{eq:direct_background_time}
\end{equation}
Accidental coincidence between physical streams is evaluated at common-time
anchors on days 0, 5, 10, 15, and 20. At each anchor, Poisson arrival times are
generated from the component rates and merged onto one time axis before the
$1\,\mu\mathrm{s}$ grouping, detector response, active anticoincidence, and
Compton topology/aperture selections. The calculation defines a background
correction and a focused-signal survival fraction,
\begin{equation}
 \widehat\rho_{g,r}=
 \frac{n^{\mathrm{tl}}_{g,r}/T^{\mathrm{rep}}_{g,r}}
 {B^{\mathrm{dir}}_g(t_r)},\qquad
 \widehat\eta_{g,r}=
 \frac{n^{\mathrm{pass}}_{g,r}}{N^{\mathrm{probe}}_{g,r}}.
 \label{eq:anchor_time_factors}
\end{equation}
$\widehat\rho$ converts the direct background to the final common-time rate.
It can be slightly larger than unity when accidental energy summation moves an
event into the science window and is therefore not a survival probability.
$\widehat\eta$ is the conditional survival fraction of a weak focused event.
The two factors are interpolated separately between adjacent anchors. The
equivalent exposure per anchor is $2\times10^4$ s for mass model A and
$2\times10^5$ s for mass model B.

Figure~\ref{fig:common_time_anchors} checks the event-time layer against the
uncorrected analytic fold before applying $\rho_g(t)$. For each mass model, a
separate Poisson replay is generated on days 0, 5, 10, 15, and 20 and passed
through the common-time grouping and event selections. The replay rates are
then compared with $B_g^{\mathrm{dir}}(t)$ from
Equation~\eqref{eq:direct_background_time}, evaluated on all 81 nodes. The
largest replay-only differences are $1.67\sigma$ for mass model A and
$1.09\sigma$ for mass model B; no anchor differs from the analytic direct-rate
curve by $2\sigma$. This agreement supports use of the analytic fold as the
time-dependent rate baseline within the sampled trajectory. The comparison is
nevertheless a closure test, not five independent particle transports: the
component rates supplied to each replay are themselves evaluated from
Equations~\eqref{eq:prompt_source_response}--\eqref{eq:delayed_selected}.

\begin{figure}[!ht]
\centering
\includegraphics[width=\linewidth]{figures/section4_current/fig04_time_fold_validation.pdf}
\caption{Event-level closure of the mission-time background fold. (a) The
solid and dashed curves are the 81-node analytic direct science-window
background rates obtained from the day-15 component responses, family-resolved
environmental scaling, and nuclide-resolved activation evolution. Shaded bands
show the finite-transport $1\sigma$ template uncertainties; because the same
transport templates are reused at all nodes, these uncertainties are
correlated in time. Open symbols show separate common-time Poisson replays on
days 0, 5, 10, 15, and 20. Their error bars are $\sqrt{n}/T$, with
$T=2\times10^4$ s for mass model A and $2\times10^5$ s for mass model B.
(b) Replay-to-analytic rate ratios with replay-only statistical errors; the
horizontal line marks unity. The transport bands and replay errors are shown
separately and are not combined. These replays test rate mapping, common-time
grouping, and event-selection closure for the supplied component rates; they
are not independent full transports at five environments.}
\label{fig:common_time_anchors}
\end{figure}

The final time-dependent rates are
\begin{equation}
 S_g(t)=S^{\mathrm{dir}}_g(t)\eta_g(t),\qquad
 B_g(t)=B^{\mathrm{dir}}_g(t)\rho_g(t).
 \label{eq:final_time_rates}
\end{equation}
Trapezoidal integration over the 81 nodes then gives the cumulative counts at
exposure $T_n$,
\begin{equation}
 \begin{aligned}
 N_s(T_n)&=\sum_{\ell=1}^{n}\frac{S^{\mathrm{dir}}_g(t_{\ell-1})\eta_g(t_{\ell-1})+S^{\mathrm{dir}}_g(t_\ell)\eta_g(t_\ell)}{2}\,\Delta t_\ell,\\
 N_b(T_n)&=\sum_{\ell=1}^{n}\frac{B^{\mathrm{dir}}_g(t_{\ell-1})\rho_g(t_{\ell-1})+B^{\mathrm{dir}}_g(t_\ell)\rho_g(t_\ell)}{2}\,\Delta t_\ell.
 \end{aligned}
 \label{eq:cumulative_counts}
\end{equation}
The resulting $N_s(T)$ and $N_b(T)$ are the inputs to the sensitivity definition
in the next subsection. Their cumulative realization is reported in
Section~\ref{subsec:mission_fold_validation}.

\subsection{Mission significance and minimum resolvable line flux}
\label{sec:mission_sensitivity}

For cumulative signal and background counts, we define the counting
significance as
\begin{equation}
 Z(T)=\frac{N_s(T)}{\sqrt{N_b(T)}}.
 \label{eq:mission_significance}
\end{equation}
Under the high-count, weak-signal conditions considered here, this is the
large-sample limit of the Poisson counting likelihood for $N_s\ll N_b$
\cite{Cowan2011}.

The focused-signal count is linear in line flux, whereas the predicted
background is independent of a weak source. The $q\sigma$ minimum resolvable
top-of-atmosphere line flux at exposure $T$ is therefore
\begin{equation}
 F_{q\sigma}(T)=\fref\frac{q}{Z(T;\fref)},
 \qquad q\in\{3,5\}.
 \label{eq:minimum_flux}
\end{equation}
The time required to reach a specified significance is the first crossing of
$Z(T)=q$.

The same transport events are reused at all 81 mission nodes, so their
statistical errors are correlated across time. We first accumulate each event's
weighted contribution over all nodes and then calculate the statistical
uncertainty of the mission-integrated counts. Statistical errors from the five
time anchors, effective area, and signal probes are propagated into the time
curves and flux threshold. Threshold-crossing times are linearly interpolated
between adjacent nodes.

\section{Results}
\label{sec:results}

This section summarizes the results within the simulation framework described above. We first establish the background and response of mass model A and trace the selected events to the materials, parent nuclides, and spatial regions that limit its performance. We then test the complete side-chimney configuration of mass model B and compare the two designs on the common mission time axis. The energy response, science window, arrival-time grouping, and Compton topology/aperture criterion are the same for both models; the active-shield channels are part of each complete design. \revblue{For diagnostic presentation, the same eight-family prompt branch is separated below into atmospheric $\gamma$ rays and the aggregate of the seven non-$\gamma$ families. The simulated streams retain the particle-family physical-rate normalization defined above: each event carries its corresponding rate weight, selected component rates are sums over the weights of accepted events, and the common-time replay draws Poisson arrivals from those component rates.}

\subsection{Baseline background and selection performance of mass model A}
\label{subsec:model_a_baseline}
\label{subsec:reference_background_results}

Mass model A defines the reference performance. Table~\ref{tab:phase2_cutflow} gives the day-15 direct, no-coincidence cut flow in the $[510.58,511.42)\keV$ science window. The active-BGO gate reduces the total background from 3,086 records and $5.43630\cps$ to 443 records and $6.07099\times10^{-2}\cps$, removing 98.883\% of the pre-veto rate. The \revblue{seven non-$\gamma$ prompt particle families} fall from $2.91217\cps$ to a zero central value, the \revblue{atmospheric-$\gamma$ component} falls from $2.35667\cps$ to $1.78923\times10^{-2}\cps$, and delayed activation falls from $0.167460\cps$ to $4.28176\times10^{-2}\cps$.

After the Compton topology/aperture selection, 400 background records remain at $5.30571\times10^{-2}\cps$, or 87.394\% of the active-veto rate. All 28,612 isolated focused-signal records in the measured window pass the active gate, and 28,006 survive the final selection. The corresponding effective area is $15.12324\pm0.04492\,\mathrm{cm^2}$, equal to 97.882\% of the measured-window signal area. Active anticoincidence supplies most of the rate reduction, while the topology selection removes a smaller residual and retains most of the focused signal.

\begin{table}[H]
\centering
\caption{Mass model A science-window cut flow. Background entries are day-15 direct, no-coincidence record counts and rates; signal entries give record counts and effective area for the common 37,194-photon focal-plane sample.}
\label{tab:phase2_cutflow}
\footnotesize
\begin{tabular}{lccc}
\toprule
Stream & Pre-veto & After active veto & Final topology/aperture \\
\midrule
\revblue{Seven non-$\gamma$ prompt particle families} & 869 / $2.91217\cps$ & 0 / 0 & 0 / 0 \\
\revblue{Atmospheric $\gamma$} & 922 / $2.35667\cps$ & 7 / $1.78923\times10^{-2}\cps$ & 6 / $1.53363\times10^{-2}\cps$ \\
Delayed activation & 1295 / $0.167460\cps$ & 436 / $4.28176\times10^{-2}\cps$ & 394 / $3.77209\times10^{-2}\cps$ \\
Total background & 3086 / $5.43630\cps$ & 443 / $6.07099\times10^{-2}\cps$ & 400 / $5.30571\times10^{-2}\cps$ \\
Focused signal & 28612 / $15.45048\,\mathrm{cm^2}$ & 28612 / $15.45048\,\mathrm{cm^2}$ & 28006 / $15.12324\,\mathrm{cm^2}$ \\
\bottomrule
\end{tabular}
\end{table}

The nonzero components of the final direct background are delayed activation
and \revblue{atmospheric $\gamma$ rays}
(Table~\ref{tab:reference_background_budget}). Delayed activation contributes
$3.77209\times10^{-2}\cps$ (71.095\%), and the \revblue{atmospheric-$\gamma$ component} contributes
$1.53363\times10^{-2}\cps$ (28.905\%). Because the event weights differ among
components, record count alone does not measure statistical support; the
400-record total has $N_{\mathrm{eff}}=46.77$.

\begin{table}[H]
\centering
\caption{Day-15 final direct background of mass model A. Standard errors and effective sample sizes are calculated from $\sum_i w_i^2$ for unequal event weights.}
\label{tab:reference_background_budget}
\footnotesize
\begin{tabular}{lrrrrr}
\toprule
Component & Records & Rate / $\cps$ & Standard error / $\cps$ & $N_{\mathrm{eff}}$ & Total / \% \\
\midrule
Delayed activation & 394 & $3.77209\times10^{-2}$ & $4.58152\times10^{-3}$ & 67.79 & 71.095 \\
\revblue{Atmospheric $\gamma$} & 6 & $1.53363\times10^{-2}$ & $6.26100\times10^{-3}$ & 6.00 & 28.905 \\
\midrule
Total & 400 & $5.30571\times10^{-2}$ & $7.75825\times10^{-3}$ & 46.77 & 100.000 \\
\bottomrule
\end{tabular}
\end{table}

\subsection{Background origins and spatial coupling in mass model A}
\label{subsec:mass_model_a_origins}

\revmark{19}
\revmark{20}
\begingroup
\color{revisionblue}
The preceding subsection establishes the background level of mass model A. To identify its physical origins, we traced all final events by initiating particle family. The delayed-activation component was further decomposed by parent-production material, source-volume/parent-nuclide pair, and spatial region relative to the TES.

Proton-induced activation is the largest single contribution, with 35 records and $2.48463\times10^{-2}\cps$, or 46.829\% of the total direct background. The remaining delayed terms are $7.49239\times10^{-3}\cps$ from $\alpha$-induced activation, $3.60054\times10^{-3}\cps$ from neutron-induced activation, $1.68617\times10^{-3}\cps$ from $\gamma$-induced activation, $7.02083\times10^{-5}\cps$ from $e^+$-induced activation, and $2.52755\times10^{-5}\cps$ from $e^-$-induced activation. Together they reproduce the final delayed rate of $3.77209\times10^{-2}\cps$.

In the material decomposition, copper accounts for 70.87\% of the final delayed rate. The three leading source-volume/nuclide pairs are $^{62}$Cu produced in the TES substrate-support copper panel, the 50 mK mixing-chamber copper cold plate, and the TES-stack copper heat-sink ring. They are followed by $^{199}$Po produced in the Bi passive-shield layer and $^{64}$Cu produced in the same 50 mK cold plate and TES substrate-support copper panel. This ordering includes particle transport, detector response, active anticoincidence, and Compton selection.

Grouping the parent-production sites by transport component identifies the 50 mK mixing-chamber copper cold plate as the largest source, at 28.59\% of the final delayed rate. The two TES substrate-support copper panels (L2 and L4) contribute 27.73\%, followed by the TES-stack copper heat-sink ring at 10.18\% and the near-field Bi passive-shield layer at 7.49\%. Aluminium cryostat and shield components contribute 15.35\% in total; the remaining dilution-refrigerator/cold-stage hardware, the tungsten detector-bay bottom plate, and the BPE plus residual BGO shielding contribute 5.26\%, 3.51\%, and 1.90\%, respectively. The first three copper component groups together account for 66.49\% of the delayed rate, locating the dominant burden in the 50 mK cold plate, TES substrate supports, and TES heat sink.

Figure~\ref{fig:activation_origin_donuts} projects the same final selected delayed-activation events onto their parent-production components and initiating primary-particle families. The component projection resolves the cold plate, TES supports, and heat sink, while the particle-family projection identifies protons and $\alpha$ particles as the main drivers of this component.

\begin{figure}[!ht]
\centering
\includegraphics[width=0.96\linewidth]{figures/section4_current/fig_activation_origin_donuts_model_a.pdf}
\caption{\revblue{Origins of the final selected day-15 delayed-activation background in mass model A. (a) Component in which the parent nuclide was produced. The TES substrate-support copper category combines the L2 and L4 panels; the aluminium cryostat/shield category combines the temperature-stage shield caps, inner shield cylinder, vacuum-jacket cap, and 50 mK aluminium-can bottom cap; ``other DR/cold-stage hardware'' contains the remaining copper DR/cold-stage parts and the silver-sinter proxy. (b) Initiating primary-particle family, with $e^+$ and $e^-$ combined as $e^{\pm}$. Sector areas and legend percentages are normalized to the final delayed-activation rate, $3.77209\times10^{-2}\cps$.}}
\label{fig:activation_origin_donuts}
\end{figure}

\revmark{22}Surrounding a main detector with active anticoincidence counters is an established low-background strategy in hard-X-ray and soft-$\gamma$-ray instruments; the Suzaku/HXD, for example, used thick BGO guard counters around its main detector units \cite{Takahashi2007HXD}. In the present instrument, placing conventional BGO directly at the mK cold end would add the scintillator, photosensors, mechanical supports, and cabling to the cold assembly, increasing both the thermal load and integration complexity. BGO light yield and pulse-decay time are also temperature dependent \cite{Verdier2011BGO}. Mass model A therefore uses a Bi passive-shield layer near the TES while retaining the active BGO outside the ultracold assembly. Tungsten is a representative high-$Z$ material for passive shielding in the immediate vicinity of high-energy radiation detectors and is therefore used here as the engineering benchmark. At an areal density adjusted to provide the same attenuation at 511 keV, the NIST XCOM mass attenuation coefficients give a lower pair-production optical depth in Bi than in W for several-MeV atmospheric $\gamma$ rays; for the $4.403906\MeV$ primary traced below, the Bi value is approximately 8\% lower \cite{NISTXCOM}.

MeV atmospheric $\gamma$ rays reaching the focal-plane region can still undergo electron--positron pair production in the nearby passive shield and generate secondary annihilation photons. Event-level process tracing identified a complete sequence: a $4.403906\MeV$ atmospheric $\gamma$ ray underwent pair production in the 4.796-mm-thick Bi passive-shield layer. After the positron annihilated in the same layer, one $510.999\keV$ photon reached the TES and underwent Compton scattering followed by photoabsorption in two pixels, while the active shield recorded zero deposited energy.

This source diagnosis directly motivates mass model B: the TES sensitive region is moved outside the direct projection of the refrigerator cold plates, and local active BGO is reorganized around the lateral focal plane to reduce both the geometric coupling to cold-stage activation and the production of secondary 511-keV photons by high-energy $\gamma$ rays in near-field passive shielding.
\endgroup

\FloatBarrier

\subsection{From the source diagnosis to the complete mass model B configuration}
\label{subsec:optimized_shield_results}

\revmark{21}\revblue{Mass model B places the six-layer TES focal plane laterally inside a side chimney rather than directly beneath the staged cold structures, while retaining the copper thermal-support frame of each TES layer. Heat from the layers is collected by the rear copper heat-sink assembly and conducted through a bent copper cold finger, which passes through a 1.50-cm-diameter rear cold port and connects to the 50 mK mixing-chamber cold plate. Except for this thermal-service port, the cold-stage-facing rear face of the focal plane is covered as fully as practicable by an active BGO rear annulus. Together with the chimney-side shield and front optical annulus, it tags 511-keV photons entering the rear of the TES from the cold-stage direction, including annihilation photons associated with $\beta^+$ decays in activated cold-stage copper. Estimated from the nominal geometry volumes using the same BGO density, the BGO mass decreases from approximately 298 kg in mass model A to approximately 40 kg in mass model B, a reduction of approximately 258 kg ($87\%$). Events retain the common $1\,\mu\mathrm{s}$ time grouping and pass the active gate when the summed BGO deposit in the group is below 50 keV.}

\revblue{Figure~\ref{fig:optimized_shield_background} compares the two cold-end
structures in the same instrument coordinates and at the same scale. The
detector-neighbourhood panels emphasize the six-layer TES, layer-by-layer copper
thermal supports, bent copper cold finger, rear cold port, and surrounding active BGO.}

\begin{figure}[!ht]
\centering
\includegraphics[width=\linewidth]{figures/section4_current/fig09_geometry_response_detailed.pdf}
\caption{\revblue{Dimensionally aligned $x'$--$z'$ sections of the cold end and focal plane, overlaid with the parent-production positions of final delayed events. Panels (a) and (b) show mass models A and B over identical coordinate ranges and at the same scale; panels (c) and (d) show their detector neighbourhoods at a common physical scale. The drawings identify the staged copper cold plates, six-layer TES, active BGO, tungsten aperture/frame, model-B chimney and copper cold finger, and the model-A TES substrate-support copper panel, TES-stack copper heat-sink ring, and 50 mK MXC-stage copper cold plate. Each symbol marks the production position associated with a day-15 delayed event that survives the final selection; colour identifies the source material, and marker area follows the selected weighted rate on a common scale subject to a minimum display size. Sections are drawn to the transport-geometry dimensions; minor CSG apertures and service components are omitted for clarity.}}
\label{fig:optimized_shield_background}
\end{figure}

\revblue{The two mass models use the same source fields, TES response, science window, common-time grouping, and Compton topology/aperture selection. Their day-15 backgrounds, focused responses, and mission performance are compared within this common simulation framework.}

\FloatBarrier

\subsection{Full-chain background and source redistribution in mass model B}
\label{subsec:optimized_fullchain}

Mass model B uses the same 420 eV FWHM TES response, $[510.58,511.42)\keV$ science window, and Compton topology/aperture test as mass model A. Table~\ref{tab:optimized_cutflow} gives its day-15 direct, no-coincidence cut flow.

Active BGO reduces the total background from 7,684 records and $2.38372\cps$ to 137 records and $1.11379\times10^{-2}\cps$, or 0.4672\% of the pre-veto rate. All 28,434 isolated focused-signal records in the measured window pass the active gate. The final Compton topology/aperture selection leaves 126 background records and retains 96.492\% of the active-veto weighted rate. It retains 27,936 focused records, corresponding to $15.08544\pm0.04503\,\mathrm{cm^2}$ and 98.249\% of the measured-window signal area. \revblue{Under the common on-axis signal sample and selection framework, the final effective area of mass model B is 99.75\% of that of mass model A.}

\begin{table}[H]
\centering
\caption{Mass model B response-convolved science-window cut flow. Background entries are day-15 direct, no-coincidence record counts and rates; signal entries give record counts and effective area.}
\label{tab:optimized_cutflow}
\footnotesize
\begin{tabular}{lccc}
\toprule
Stream & Pre-veto & After active BGO & Final topology/aperture \\
\midrule
\revblue{Seven non-$\gamma$ prompt particle families} & 1448 / $0.878692\cps$ & 3 / $1.81131\times10^{-3}\cps$ & 3 / $1.81131\times10^{-3}\cps$ \\
\revblue{Atmospheric $\gamma$} & 2987 / $1.37944\cps$ & 12 / $5.35535\times10^{-3}\cps$ & 12 / $5.35535\times10^{-3}\cps$ \\
Delayed activation & 3249 / $0.125586\cps$ & 122 / $3.97120\times10^{-3}\cps$ & 111 / $3.58054\times10^{-3}\cps$ \\
Total background & 7684 / $2.38372\cps$ & 137 / $1.11379\times10^{-2}\cps$ & 126 / $1.07472\times10^{-2}\cps$ \\
Focused signal & 28434 / $15.35436\,\mathrm{cm^2}$ & 28434 / $15.35436\,\mathrm{cm^2}$ & 27936 / $15.08544\,\mathrm{cm^2}$ \\
\bottomrule
\end{tabular}
\end{table}

The final model-B direct background contains $1.81131\times10^{-3}\cps$ from
the \revblue{seven non-$\gamma$ prompt particle families}, $5.35535\times10^{-3}\cps$ from \revblue{atmospheric
$\gamma$ rays}, and $3.58054\times10^{-3}\cps$ from delayed
activation. Their shares are 16.854\%, 49.830\%, and 33.316\%, respectively.
Relative to mass model A, the delayed and \revblue{atmospheric-$\gamma$ rates} fall to
9.492\% and 34.920\%, and the total rate is $1.07472\times10^{-2}\cps$, or
20.256\% of the model-A value, a factor of 4.94 lower. \revblue{Table~\ref{tab:optimized_component_precision}
lists the component rates and finite-transport standard errors, and
Figure~\ref{fig:background_composition} compares the final fractional
compositions of the two models using the same categories.}

\begin{table}[!ht]
\centering
\caption{Day-15 final direct-background composition and statistical uncertainty for mass model B. Standard errors and effective sample sizes are calculated from $\sum_i w_i^2$ for unequal event weights.}
\label{tab:optimized_component_precision}
\scriptsize
\resizebox{\linewidth}{!}{%
\begin{tabular}{lrrrrr}
\toprule
Component & Records & Rate / $\cps$ & Standard error / $\cps$ & $N_{\mathrm{eff}}$ & Total / \% \\
\midrule
\revblue{Seven non-$\gamma$ prompt particle families} & 3 & $1.81131\times10^{-3}$ & $1.04576\times10^{-3}$ & 3.00 & 16.854 \\
\revblue{Atmospheric $\gamma$} & 12 & $5.35535\times10^{-3}$ & $1.55631\times10^{-3}$ & 11.84 & 49.830 \\
Delayed activation & 111 & $3.58054\times10^{-3}$ & $5.36971\times10^{-4}$ & 44.46 & 33.316 \\
\midrule
Total & 126 & $1.07472\times10^{-2}$ & $1.95040\times10^{-3}$ & 30.36 & 100.000 \\
\bottomrule
\end{tabular}}
\end{table}

Copper-derived events contribute $(2.26708\pm0.42826)\times10^{-3}\cps$, or
63.32\% of the final model-B delayed rate. The contribution from the DR/mixing chamber and staged cold plates changes from
$(1.21398\pm0.26869)\times10^{-2}\cps$ in mass model A to
$(1.40604\pm1.09289)\times10^{-4}\cps$ in mass model B. The
total delayed rate changes from $(3.77209\pm0.45815)\times10^{-2}$ to
$(3.58054\pm0.53697)\times10^{-3}\cps$.

Figure~\ref{fig:activation_origin_donuts_b} resolves the 111 final delayed-activation
origins in mass model B. The TES substrate-support Cu panels and TES Cu
heat-sink ring contribute 26.05\% and 25.73\%, respectively, or 51.78\%
together. The W focal-plane frame and the Al chimney shells, service lid, and annuli each
contribute 16.42\%. The TES bottom Cu cold-plate spoke, 50~mK mixing-chamber
Cu plate, and cold-link hardware contribute 5.72\%, 3.93\%, and 3.78\%,
respectively, while the TES-pixel Ta volumes and residual BGO side shield each
contribute about 1\%. The initiating primaries are dominated by protons
(57.89\%), followed by neutrons (20.81\%) and $\alpha$ particles (18.39\%);
$\gamma$ rays and leptons contribute 2.83\% and 0.09\%, respectively. Thus,
the residual delayed background in mass model B is concentrated in the TES
support panels and heat-sink ring, whereas the 50~mK mixing-chamber Cu plate
accounts for 3.93\%.

\begin{figure}[!htbp]
\centering
\includegraphics[width=0.96\linewidth]{figures/section4_current/fig_activation_origin_donuts_model_b.pdf}
\caption{Origins of the final day-15 delayed-activation background in mass model B. (a) Component in which the parent nuclide was produced. (b) Primary particle that initiated its production. Four Cu substrate-support panels (L2/L4/L5) are grouped, as are the Al chimney/port/lid/annulus volumes and the Cu-finger/SS-rod cold link; Ta pixels and BGO remain separate. Leptons combine $e^-$, $e^+$, and $\mu^-$. Each of the 111 records was reweighted by primary family to the current formal rate before aggregation. Both panels use $3.58054\times10^{-3}\cps$ as their denominator.}
\label{fig:activation_origin_donuts_b}
\end{figure}

\begin{figure}[H]
\centering
\includegraphics[width=0.92\linewidth]{figures/section4_current/fig09_background_composition.pdf}
\caption{Composition of the final direct background in the day-15 $510.58$--$511.42\keV$ science window after all selections. Sector area gives each physical component's fraction of the corresponding model total, shown at the centre; the same component colours are used in both panels.}
\label{fig:background_composition}
\end{figure}

Figure~\ref{fig:day15_selection_spectra} shows the day-15 spectra before active
anticoincidence, after the active gate, and after the final Compton selection.
The 480--550 keV band provides a selection diagnostic independent of the
narrow-window integral. In mass model A, its weighted rate changes from
$10.7984\cps$ to $0.140797\cps$ and then $0.121285\cps$; the corresponding
model-B rates are $4.92157$, $0.0735321$, and $0.0655197\cps$. The same ordering
appears in the science-window cut flows: active anticoincidence supplies the
dominant suppression, and the Compton topology/aperture selection makes a
smaller correction to the residual.

\begin{figure}[!ht]
\centering
\includegraphics[width=\linewidth]{figures/section4_current/fig05_day15_selection_spectra_current.pdf}
\caption{\revblue{Day-15 detector-side differential count-rate spectra through the event-selection chain. Panels (a) and (b) show 480--550 keV for mass models A and B; four adjacent 0.25 keV analysis bins are summed into each 1 keV display bin. Panels (c) and (d) show 508--514 keV at the native 0.25 keV binning; the shaded interval and dotted boundaries mark the $[510.58,511.42)\keV$ science window. Curves show the direct, no-coincidence background before active anticoincidence, after active anticoincidence, and after the final Compton topology/aperture selection. The translucent band in the 480--550 keV panels and the bin-wise error bars in the zoom panels show finite-transport $1\sigma$ uncertainties from $\sqrt{\sum_i w_i^2}$. Zero-weight bins are left blank, without smoothing or pseudocounts. The ordinate is the detector differential count rate after transport and TES-response convolution.}}
\label{fig:day15_selection_spectra}
\end{figure}

Both delayed activation and the \revblue{atmospheric-$\gamma$ component} have
lower rates in mass model B than in mass model A. The dominant component
changes from delayed activation in mass model A to \revblue{atmospheric
$\gamma$ rays} in mass model B.

\FloatBarrier

\subsection{Common-time fold and cumulative counts}
\label{subsec:mission_fold_validation}

At the day-0, 5, 10, 15, and 20 anchors, Poisson arrivals are regenerated from the physical component rates and passed through the common time grouping, detector response, active anticoincidence, and Compton selection. The ratio of the common-time final rate to the direct rate defines $\widehat\rho_g$, and focused-signal replay gives the accidental survival $\widehat\eta_g$; both are interpolated to the 81 mission nodes. The model-A $\widehat\rho_g$ values span 0.9307--1.0060, and the model-B values span 0.9901--1.0293 (Figure~\ref{fig:common_time_anchors}).

As shown in Figure~\ref{fig:common_time_anchors}, all five event-level replay
rates are statistically compatible with the corresponding analytic direct-rate
curve at the replay-only $2\sigma$ level. The comparison establishes numerical
closure of the common-time implementation and quantifies the correction
$\widehat\rho_g$; it does not substitute for off-reference particle transport.

The focused-signal $\widehat\eta_g$ values span 0.96810--0.97039 for mass model A and 0.99547--0.99580 for mass model B; their day-15 values are 0.9682365 and 0.995472. The higher model-B survival follows from lower coincidence occupancy rather than a different signal selection. The interpolated $\rho_g$ and $\eta_g$ values are combined with the same trajectory transmission and component responses and integrated over 81 nodes at 0.25 d spacing.

\revblue{Combining the day-15 conditional survival fractions with the isolated-event final effective areas and the atmospheric transmission at that node gives final signal rates of $2.29144\times10^{-3}$ and $2.35000\times10^{-3}\cps$ for mass models A and B, respectively, at the reference flux defined in Section~\ref{sec:requirements}. Mass model B has a slightly smaller isolated-event effective area but a smaller accidental-coincidence loss, and therefore a slightly higher final reference signal rate.}

At 20 d, mass model A accumulates 3928.83 signal counts at the reference flux
and 87632.34 background counts; mass model B accumulates 4028.15 and 18568.89.
The \revblue{atmospheric $\gamma$} and remaining-background totals are
25760.39 and 61871.95 counts for mass model A and 9359.44 and 9209.45 counts for
mass model B. The model-B cumulative background is 21.19\% of the model-A
value.

\subsection{Mission significance and minimum resolvable line flux}
\label{subsec:final_mission_performance}

Figure~\ref{fig:optimized_mission_significance} and
Table~\ref{tab:primary_sensitivity} summarize the mission performance. At
$\fref$, the 20 d counting significances are 13.272 for mass model A and
29.561 for mass model B. The corresponding $3\sigma$ minimum resolvable
top-of-atmosphere fluxes are $(5.425\pm0.403)\times10^{-5}$ and
$(2.436\pm0.223)\times10^{-5}\phcms$.

\begin{figure}[!ht]
\centering
\includegraphics[width=0.98\linewidth]{figures/section4_reference_flux_20260830/fig10_mission_performance_gaussian.pdf}
\caption{Conditional 81-node mission performance of mass models A and B. (a)
Counting significance at the top-of-atmosphere reference flux
$\fref=2.4\times10^{-4}\phcms$; horizontal lines mark $3\sigma$ and
$5\sigma$. (b) Corresponding $3\sigma$ minimum resolvable top-of-atmosphere
flux. Both designs use the same source, detector response, science window,
timing, topology/aperture test, and atmospheric-transmission conventions; their
entrances, local tungsten structures, and geometry-specific active channels
differ.}
\label{fig:optimized_mission_significance}
\end{figure}

\begin{table}[!ht]
\centering
\caption{Conditional mission projection for mass models A and B.}
\label{tab:primary_sensitivity}
\footnotesize
\resizebox{\linewidth}{!}{%
\begin{tabular}{lrr}
\toprule
Quantity & Mass model A & Mass model B \\
\midrule
Selected effective area / $\mathrm{cm^2}$ & $15.12324\pm0.04492$ & $15.08544\pm0.04503$ \\
Day-15 direct background / $\cps$ & $5.30571\times10^{-2}$ & $1.07472\times10^{-2}$ \\
Day-15 signal at $\fref$ / $\cps$ & $2.29144\times10^{-3}$ & $2.35000\times10^{-3}$ \\
20 d background counts & 87632.34 & 18568.89 \\
20 d signal counts at $\fref$ & 3928.83 & 4028.15 \\
$Z_{20\mathrm d}=N_s/\sqrt{N_b}$ & 13.272 & 29.561 \\
$T_{3\sigma}$ at $\fref$ / d & 0.669 & 0.198 \\
$T_{5\sigma}$ at $\fref$ / d & 1.913 & 0.452 \\
$F_{3\sigma}(20\,\mathrm d)$ / $\phcms$ &
$(5.425\pm0.403)\times10^{-5}$ &
$(2.436\pm0.223)\times10^{-5}$ \\
$F_{5\sigma}(20\,\mathrm d)$ / $\phcms$ &
$(9.042\pm0.671)\times10^{-5}$ &
$(4.059\pm0.372)\times10^{-5}$ \\
\bottomrule
\end{tabular}}
\end{table}

\FloatBarrier

The flux-threshold ratio separates into cumulative-background and cumulative-signal-kernel terms. The 20 d model-B background is 21.19\% of the model-A value, while its cumulative signal kernel is 1.025278 times the model-A value:
\begin{equation}
 \frac{F_{\min,A}}{F_{\min,B}}
 =\sqrt{\frac{N_{b,A}}{N_{b,B}}}\frac{K_B^{20\mathrm d}}{K_A^{20\mathrm d}}
 =2.172397\times1.025278
 =2.227311 .
 \label{eq:fmin_ratio_decomposition}
\end{equation}
Here $K_g^{20\mathrm d}$ is the cumulative signal kernel integrated over the 81 nodes. Equation~\eqref{eq:fmin_ratio_decomposition} shows that most of the 2.227-fold threshold difference follows from the square root of the background ratio, with a smaller correction from the cumulative signal kernel.

Mass model B retains 99.75\% of the focused effective area of mass model A and reduces the 20 d cumulative background and $3\sigma$ flux threshold to 21.19\% and 44.90\%, respectively.

\FloatBarrier

\section{Response screening in orbital and lunar environments}
\label{sec:environment_screening}

The balloon calculation establishes the background and sensitivity of mass
model B at 38\,km. To identify which radiation fields would limit the same
detector response in other mission environments, we traced the selected events
back to their primary particle types and energies and reweighted them with
source spectra for LEO, the lunar surface, and Sun--Earth L2. The calculation
therefore shows how a change in radiation environment acts on the established
TES response and helps set priorities for detailed follow-up simulations.

\subsection{Spectral reweighting}

We considered four representative source fields: quiet, near-equatorial LEO at
530\,km outside the SAA; an exposed lunar surface; and two Sun--Earth L2
conditions representing solar maximum and solar minimum. The LEO spectra use
the components adopted in low-Earth-orbit gamma-ray background studies
\cite{Cumani2019LEO}. The lunar field combines the upward REDMoon spectrum for
Apollo-17 regolith with the cosmic gamma-ray sky visible from the surface
\cite{Dobynde2021REDMoon}. The L2 fields contain gamma rays, protons, alpha
particles, electrons, and positrons, with GCR spectra appropriate to the two
solar-modulation conditions
\cite{Lotti2021Athena,Usoskin2005,Kuznetsov2017}.

For the \revblue{prompt-gamma background}, each selected event retains the energy of
its primary photon. A selected delayed event is linked instead to the
activation-production energy distribution for the same incident particle,
parent nuclide, and production region. At those primary energies, each balloon
event weight is multiplied by the ratio of the target spectrum to the balloon
spectrum. The two components are then normalized to the formal day-15 rates of
mass model B: $5.355349\times10^{-3}\cps$ for \revblue{prompt gamma rays} and
$3.580538\times10^{-3}\cps$ for delayed activation. The principal response
bands shown in Figure~\ref{fig:environment_transfer} are
$0.598$--$29.6$\,MeV for selected prompt gamma rays and
$1.37$--$48.4$\,GeV for proton-induced delayed background.

The combined contribution from the seven non-gamma prompt families is 16.854\%
of the balloon background, but the present transfer response has no independent
primary-energy kernel for the prompt-positron component. In the environmental
conversion, this contribution follows the relative change of the combined
\revblue{prompt-gamma} and delayed-activation rates; no separate particle-family rate
is quoted for the target environments.

All scenarios use the same 20\,d exposure and the effective area of mass model
B. The screening flux is anchored to the balloon result of
$2.435681\times10^{-5}\phcms$, scaled with the square root of the reweighted
background and corrected for the cumulative signal loss of the balloon
trajectory relative to an unattenuated, continuously on-source observation.
The balloon value is quoted at the top of the atmosphere, whereas the other
values refer to the local instrument incident plane.

\begin{figure}[!ht]
\centering
\includegraphics[width=\linewidth]{figures/environment_screening_20260830/fig11_environment_transfer_en.pdf}
\caption{Environmental spectral reweighting for mass model B. (a) Gamma-ray
source spectra; the shaded band is the principal primary-energy range of the
selected prompt-gamma events ($0.598$--$29.6$\,MeV). (b) Proton spectra; the
shaded band is the principal response range of the proton-induced delayed
background ($1.37$--$48.4$\,GeV). (c) The 20\,d, $3\sigma$ screening flux
anchored to the balloon result. The balloon value is defined at the top of the
atmosphere and the others at the local instrument incident plane.}
\label{fig:environment_transfer}
\end{figure}

\subsection{Screening results for LEO, the lunar surface, and L2}

\begin{table}[!ht]
\centering
\caption{Spectral-reweighting screen for mass model B in different radiation
environments. The third column gives the combined \revblue{prompt-gamma} and
delayed-activation rate relative to the balloon reference; the last column is
the $3\sigma$ screening flux for the same 20\,d exposure.}
\label{tab:environment_screening}
\small
\resizebox{\linewidth}{!}{%
\begin{tabular}{llrr}
\toprule
Scenario & Source-field proxy & Reweighted background/balloon & 20\,d $3\sigma$ screening flux / $\phcms$ \\
\midrule
38 km balloon & formal 81-node trajectory & 1.0000 & $2.4357\times10^{-5}$ \\
530 km LEO & near-equatorial, quiet non-SAA & 0.5627 & $1.1764\times10^{-5}$ \\
Exposed lunar surface & REDMoon Apollo-17 regolith & 23.0228 & $7.5247\times10^{-5}$ \\
Sun--Earth L2 (2014) & solar maximum, GCR-low & 5.9460 & $3.8241\times10^{-5}$ \\
Sun--Earth L2 (2009) & solar minimum, GCR-high & 11.9177 & $5.4139\times10^{-5}$ \\
\bottomrule
\end{tabular}}
\end{table}

For quiet non-SAA LEO, the reweighted background is 0.5627 times the balloon
reference. After removing the signal loss associated with the balloon
atmosphere and observing trajectory, the 20\,d screening flux is
$1.1764\times10^{-5}\phcms$. Thus, in quiet orbital segments, the current source
field produces less background in mass model B than the balloon environment.
For a 550\,km equatorial orbit, Cumani et al. estimated that the region occupied
by trapped particles above 0.1\,MeV accounts for about 18\% of the orbital
period; science observations are normally suspended during these SAA passages
\cite{Cumani2019LEO}. A 20\,d calendar interval would therefore provide about
16.4\,d of live exposure before any post-SAA cooldown. The 20\,d value in
Table~\ref{tab:environment_screening} denotes live science exposure. The next
LEO calculation should concentrate on activation by SAA and trapped particles
and on the delayed background after an SAA passage.

For the exposed lunar surface, the reweighted background rises to 23.0228 times
the balloon reference and the screening flux becomes
$7.5247\times10^{-5}\phcms$. The increase is driven mainly by upward regolith
radiation in the response band of the selected gamma-ray events. A lunar
application therefore requires a local gamma-ray field that couples the
regolith, detector, and lander structure rather than an isolated-payload model.

At Sun--Earth L2, the result follows solar modulation. The 2014
solar-maximum/GCR-low field gives $3.8241\times10^{-5}\phcms$, whereas the 2009
solar-minimum/GCR-high field gives $5.4139\times10^{-5}\phcms$. Their ordering
tracks the proton spectra and identifies GCR-proton activation as the main L2
driver. Because the present delayed response is supported by a limited set of
activation templates, these values are most useful for setting the priority of
a dedicated proton-activation calculation.

The reweighting therefore points to a different next step in each environment:
an orbital history including SAA and trapped particles for LEO, local
regolith--lander coupling for the lunar surface, and higher-statistics GCR
proton activation at L2. Matched prompt, activation, and delayed transport in
each target environment is the next step toward mission-level performance
predictions.

\FloatBarrier

\section{Summary}
\label{sec:summary}

The day-15 selected background rate of mass model A is
$5.30571\times10^{-2}\cps$, with delayed activation as the largest component.
Selected-event provenance assigns 32.18\% of the weighted delayed rate to the
dilution-refrigerator mixing chamber and staged cold region, while copper
accounts for 70.87\% in the \revblue{material decomposition}. The leading structures are the
50 mK copper cold plate, the TES substrate-support copper panels, and the TES
copper heat-sink ring. \revblue{The science-window background is governed by the
distance, direct lines of sight, and shield coverage between activated material
and the TES.}

Mass model B places the six TES layers laterally in a BGO-surrounded chimney.
Its final day-15 background rate is $1.07472\times10^{-2}\cps$, composed of
49.830\% \revblue{atmospheric gamma rays}, 33.316\% delayed activation, and
16.854\% from the seven non-gamma prompt families. The active BGO reduces the
direct science-window rate from $2.38372$ to
$1.11379\times10^{-2}\cps$, after which the topology and aperture selection
leaves $1.07472\times10^{-2}\cps$. The selected effective area of mass model B
is 99.75\% of that of mass model A, while its final topology and aperture
selection retains 98.249\% of the focused signal in its own measured-energy
window. Its 20\,d, top-of-atmosphere $3\sigma$ minimum resolvable line flux is
$(2.436\pm0.223)\times10^{-5}\phcms$, a factor of 2.227 lower than for mass
model A.

The results lead to three design priorities. The TES should be displaced from
near-field materials in which positron-emitting radionuclides and other
activation products are generated, while direct, unshielded lines of sight from
these components to the focal plane should be minimized. The active BGO should
cover as much solid angle as the optical path permits.
Activation-prone support, heat-sink, and shielding structures near the TES
should be reduced, replaced, or relocated, with distance and shadow shielding
used to isolate unavoidable activated components. In mass model B, the TES
substrate-support copper panels and copper heat-sink ring still account for
51.78\% of the residual delayed background and are the direct targets of the
next near-field redesign.

\section{Outlook}
\label{sec:outlook}

The most immediate extension is a dedicated LEO satellite. The preliminary
screening in Section~\ref{sec:environment_screening} identifies quiet non-SAA
orbital segments as a priority for a complete mission design. Compared with a
finite-duration balloon flight, a LEO satellite avoids attenuation through the
balloon atmospheric column and can extend live exposure from months to years.
The orbit, shielding geometry, and observing schedule should be designed
together: a near-equatorial, low-inclination orbit, shutdown during SAA
passages, planned post-SAA cooldown, reduced activation near the TES, and wider
BGO solid-angle coverage are the first parameters to study. A mission model
including the spacecraft, Earth-albedo radiation, trapped-particle and SAA
histories, Earth occultation, and source visibility can then determine how far
long-duration operation lowers the minimum resolvable line flux.

A lunar instrument is a more exploratory, longer-term option. The simple
spectral reweighting in Section~\ref{sec:environment_screening} exposes the
pressure from upward regolith gamma rays, but it does not use the additional
mass and installation volume that a lunar platform may provide. If those
resources are available, the TES, cryogenic system, and lander can be optimized
in a coupled mass model. Directional lower-hemisphere shielding, active
anticoincidence below and beside the focal plane, and shadowing by the lander
could reduce direct coupling of the upward regolith field to the TES. Coupled
regolith--lander--instrument transport will be needed to determine whether such
a design can suppress the regolith-albedo background while benefiting from
unattenuated celestial photons and long cumulative exposure.

\newpage
\begin{samepage}
A further detector-level study should address neutron elastic scattering in
the cryogenic Si wafer and nearby focal-plane supports. Nuclear-recoil energy
deposited in these volumes can generate athermal and thermal phonons; if a
fraction couples to the TES, the response may appear as a non-photon pulse or
as thermal cross-talk \cite{Agnese2018NuclearRecoil,Hull2023ThermalBackground}.
Coupled neutron-transport, substrate-phonon, TES thermal-network, and measured
pulse-template studies, together with neutron-beam calibration, are needed to
determine whether such events can be reconstructed in the 511 keV science
window and whether their pulse shapes permit rejection.
\end{samepage}

\section*{Data and software availability}
Before publication, a versioned public repository or institutional archive should provide the geometry and source definitions, stochastic-analysis conventions, analysis and figure-generation scripts, LaTeX source, and derived validation products needed to reproduce the tables and figures in this paper.

\section*{Declarations}
\textbf{Funding} To be completed before journal submission.

\textbf{Competing interests} To be completed before journal submission.

\textbf{Author contributions} To be completed before journal submission.

\textbf{Ethics approval and consent to participate} Not applicable.

\begin{thebibliography}{99}

\bibitem{Prantzos2011}
N. Prantzos, C. Boehm, A.M. Bykov, R. Diehl, K. Ferriere, N. Guessoum, P. Jean, J. Knoedlseder, A. Marcowith, I.V. Moskalenko, A. Strong, G. Weidenspointner, The 511 keV emission from positron annihilation in the Galaxy, Rev. Mod. Phys. 83 (2011) 1001--1056, doi:10.1103/RevModPhys.83.1001.

\bibitem{Jean2006}
P. Jean, J. Knoedlseder, V. Lonjou, M. Allain, J.-P. Roques, G.K. Skinner, B.J. Teegarden, G. Vedrenne, P. von Ballmoos, B. Cordier, P. Caraveo, R. Diehl, P. Durouchoux, P. Leleux, R. Schanne, G. Weidenspointner, Spectral analysis of the Galactic $e^+e^-$ annihilation emission, Astron. Astrophys. 445 (2006) 579--589, doi:10.1051/0004-6361:20053765.

\bibitem{Kierans2020}
C.A. Kierans, S.E. Boggs, A. Zoglauer, A.W. Lowell, C. Sleator, J. Beechert, T.J. Brandt, P. Jean, H. Lazar, J. Roberts, T. Siegert, J.A. Tomsick, P. von Ballmoos, Detection of the 511 keV Galactic positron annihilation line with COSI, Astrophys. J. 895 (2020) 44, doi:10.3847/1538-4357/ab89a9.

\bibitem{Siegert2016AA}
T. Siegert, R. Diehl, G. Khachatryan, M.G.H. Krause, F. Guglielmetti, J. Greiner, A.W. Strong, X. Zhang, Gamma-ray spectroscopy of positron annihilation in the Milky Way, Astron. Astrophys. 586 (2016) A84, doi:10.1051/0004-6361/201527510.

\bibitem{Siegert2016V404}
T. Siegert, R. Diehl, J. Greiner, M.G.H. Krause, A.M. Beloborodov, M. Cadolle Bel, F. Guglielmetti, J. Rodriguez, A.W. Strong, X. Zhang, Positron annihilation signatures associated with the outburst of the microquasar V404 Cygni, Nature 531 (2016) 341--343, doi:10.1038/nature16978.

\bibitem{Knodlseder2005AllSky}
J. Knoedlseder, P. Jean, V. Lonjou, G. Weidenspointner, N. Guessoum, R. Diehl, W. Gillard, M.J. Harris, G.K. Skinner, P. von Ballmoos, G. Vedrenne, P. Caraveo, B. Cordier, V. Schoenfelder, B.J. Teegarden, The all-sky distribution of 511 keV electron-positron annihilation emission, Astron. Astrophys. 441 (2005) 513--532, doi:10.1051/0004-6361:20042063.

\bibitem{Yoneda2025}
H. Yoneda, T. Siegert, S. Mittal, Imaging the positron annihilation line with 20 years of INTEGRAL/SPI observations, arXiv:2509.01066, 2025.

\bibitem{Siegert2020COSI}
T. Siegert, S.E. Boggs, J.A. Tomsick, A.C. Zoglauer, C.A. Kierans, C.C. Sleator, J. Beechert, T.J. Brandt, P. Jean, H. Lazar, A.W. Lowell, J.M. Roberts, P. von Ballmoos, Imaging the 511 keV positron annihilation sky with COSI, Astrophys. J. 897 (2020) 45, doi:10.3847/1538-4357/ab9607.

\bibitem{Siegert2019Kinematics}
T. Siegert, R. Diehl, G. Khachatryan, M.G.H. Krause, F. Guglielmetti, J. Greiner, A.W. Strong, X. Zhang, Constraints on positron annihilation kinematics in the inner Galaxy, Astron. Astrophys. 627 (2019) A126, doi:10.1051/0004-6361/201833856.

\bibitem{DeCesare2011}
G. De Cesare, Searching for the 511 keV annihilation line from galactic compact objects with the IBIS gamma ray telescope, Astron. Astrophys. 531 (2011) A56, doi:10.1051/0004-6361/201116516.

\bibitem{Barriere2009Laue}
N.M. Barriere, L. Natalucci, N. Abrosimov, P. von Ballmoos, P. Bastie, P. Courtois, M. Jentschel, J. Knodlseder, J. Rousselle, P. Ubertini, Soft gamma-ray optics: new Laue lens design and performance estimates, arXiv:0910.0488, 2009.

\bibitem{Barriere2011LauePrototype}
N.M. Barriere, J.A. Tomsick, S.E. Boggs, A. Lowell, P. von Ballmoos, Developing a second generation Laue lens prototype: high reflectivity crystals and accurate assembly, arXiv:1111.6700, 2011.

\bibitem{Weidenspointner2006MAX}
G. Weidenspointner, C.B. Wunderer, N. Barriere, A. Zoglauer, P. von Ballmoos, Monte Carlo study of detector concepts for the MAX Laue Lens gamma-ray telescope, arXiv:astro-ph/0603152, 2006.

\bibitem{Shirazi2023}
F. Shirazi, E. Gau, M.A. Hossen, D. Becker, D. Schmidt, D. Swetz, D. Bennett, D. Braun, F. Kislat, J. Gard, J. Mates, J. Weber, N.R. Cavero, S. Chun, L. Lisalda, A. West, B. Dev, F. Ferrer, R. Bose, J. Ullom, H. Krawczynski, The 511-CAM Mission: A pointed 511 keV gamma-ray telescope with a focal plane detector made of stacked transition edge sensor microcalorimeter arrays, J. Astron. Telesc. Instrum. Syst. 9 (2023) 024006, doi:10.1117/1.JATIS.9.2.024006.

\bibitem{Zhang2022TESGamma}
S. Zhang, J.-K. Xia, T. Sun, et al., Transition edge sensor based detector: from X-ray to gamma-ray, arXiv:2204.00010, 2022.

\bibitem{Hu2026}
K. Hu, M. Fritts, D. Becker, D. Schmidt, F. Zhou, A. Dasgupta, F. Kislat, M. Keller, S. Chun, A.G. Detoito, E. Gau, X. Jin, D. Bennett, D. Braun, J. Gard, J.A.B. Mates, J. Nobles, D. Radomski, N. Rodriguez Cavero, R. Snodgrass, D. Swetz, J. Ullom, S. Vengalil Menon, J. Weber, K. Wimalasena, H. Krawczynski, Design of a mission to measure the shape and substructure of the 511 keV gamma-ray line from the center of the Milky Way, J. Astron. Telesc. Instrum. Syst. 12 (2026) 025002, doi:10.1117/1.JATIS.12.2.025002.

\bibitem{Gallego2025Balloon}
S. Gallego, U. Oberlack, J. Lommler, C.M. Karwin, A. Zoglauer, P. Jean, P. von Ballmoos, C. Kierans, C.C. Sleator, J.A. Tomsick, S.E. Boggs, Bottom-up background simulations of the 2016 COSI balloon flight, Astrophys. J. 986 (2025) 116, doi:10.3847/1538-4357/add6a0.

\bibitem{Tian2026DIXE}
R. Tian, J. Mao, J. Liu, H. Jin, W. Cui, Simulation of non-X-ray background for the DIffuse X-ray Explorer (DIXE) mission, arXiv:2604.13569, 2026.

\bibitem{Caputo2017}
R. Caputo, F. Kislat, J. Racusin, on behalf of the AMEGO Team, AMEGO: Simulations of the instrument performance, Proc. 35th International Cosmic Ray Conference, PoS(ICRC2017) 783 (2018), doi:10.22323/1.301.0783.

\bibitem{Gottardi2021TESReview}
L. Gottardi, K. Nagayoshi, A review of X-ray microcalorimeters based on superconducting transition edge sensors for astrophysics and particle physics, Applied Sciences 11 (2021) 3793, doi:10.3390/app11093793.

\bibitem{Yan2019TaHeatCapacity}
D. Yan, L.M. Gades, T. Guruswamy, U.M. Patel, O. Quaranta, A. Miceli, Modelling a transition-edge sensor X-ray microcalorimeter linear array for Compton profile measurements and energy dispersive diffraction, IEEE Trans. Appl. Supercond. 29(5) (2019) 1--4, doi:10.1109/TASC.2019.2906255.

\bibitem{Xie2026AlMn}
L. Xie, Y. Zhang, Z. Li, Z. Liu, S. Shu, J. Zhou, X. Li, H. Li, H. Gao, Y. Gu, X. Lu, Y. Zhao, C. Liu, Beyond one-thousandth energy resolution with an AlMn TES detector, arXiv:2602.11728, 2026.

\bibitem{Vavagiakis2017AlMnMagnetic}
E.M. Vavagiakis, S.W. Henderson, K. Zheng, H.-M. Cho, N.F. Cothard, B. Dober, S.M. Duff, P.A. Gallardo, G. Hilton, J. Hubmayr, K.D. Irwin, B.J. Koopman, D. Li, F. Nati, M.D. Niemack, C.D. Reintsema, S. Simon, J.R. Stevens, A. Suzuki, B. Westbrook, Magnetic sensitivity of AlMn TESes and shielding considerations for next-generation CMB surveys, arXiv:1710.08456, 2017.

\bibitem{NISTXCOM}
J.H. Hubbell, S.M. Seltzer, XCOM: Photon Cross Section Database, NIST Standard Reference Database 8 (XGAM), National Institute of Standards and Technology, doi:10.18434/T48G6X.

\bibitem{Sturner2003}
S.J. Sturner, C.R. Shrader, G. Weidenspointner, B.J. Teegarden, D. Attie, B. Cordier, R. Diehl, C. Ferguson, P. Jean, A. von Kienlin, P. Paul, F. Sanchez, S. Schanne, P. Sizun, G. Skinner, C.B. Wunderer, Monte Carlo simulations and generation of the SPI response, Astron. Astrophys. 411 (2003) L81--L84, doi:10.1051/0004-6361:20031171.

\bibitem{Agostinelli2003}
S. Agostinelli, J. Allison, K. Amako, et al., Geant4---a simulation toolkit, Nucl. Instrum. Methods Phys. Res. A 506 (2003) 250--303, doi:10.1016/S0168-9002(03)01368-8.

\bibitem{Allison2016}
J. Allison, K. Amako, J. Apostolakis, et al., Recent developments in Geant4, Nucl. Instrum. Methods Phys. Res. A 835 (2016) 186--225, doi:10.1016/j.nima.2016.06.125.

\bibitem{SanchezDelRio2011XOP}
M. Sanchez del Rio, R.J. Dejus, XOP v2.4: recent developments of the x-ray optics software toolkit, Proc. SPIE 8141 (2011) 814115, doi:10.1117/12.893911.

\bibitem{Guan2023BraggReflection}
F. Guan, M. Asai, D.A. Bartkoski, M. Kleckner, Z. Harel, M. Salehpour, Adding the X-ray Bragg reflection physical process in crystal to the Geant4 Monte Carlo simulation toolkit, part I: reflection from a crystal slab, Precision Radiation Oncology 7(1) (2023) 59--66.

\bibitem{Reiazi2025G4BraggReflection}
R. Reiazi, D. Lee, D. Bartkoski, M. Kleckner, M. Salehpour, G4BraggReflection for accurate modeling of Bragg reflection in perfect and mosaic crystals, J. Phys. D: Appl. Phys. 58 (2025) 295102, doi:10.1088/1361-6463/aded1d.

\bibitem{Zoglauer2006}
A. Zoglauer, R. Andritschke, F. Schopper, MEGAlib: The Medium Energy Gamma-ray Astronomy Library, New Astron. Rev. 50 (2006) 629--632, doi:10.1016/j.newar.2006.06.049.

\bibitem{Sato2015EXPACS}
T. Sato, Analytical model for estimating terrestrial cosmic ray fluxes nearly anytime and anywhere in the world: extension of PARMA/EXPACS, PLoS ONE 10 (2015) e0144679, doi:10.1371/journal.pone.0144679.

\bibitem{Sato2016PARMA4}
T. Sato, Analytical model for estimating the zenith angle dependence of terrestrial cosmic ray fluxes, PLoS ONE 11 (2016) e0160390, doi:10.1371/journal.pone.0160390.

\bibitem{Kondev2021NUBASE}
F.G. Kondev, M. Wang, W.J. Huang, S. Naimi, G. Audi, The NUBASE2020 evaluation of nuclear physics properties, Chinese Phys. C 45 (2021) 030001, doi:10.1088/1674-1137/abddae.

\bibitem{Kierans2017}
C.A. Kierans, S.E. Boggs, J.-L. Chiu, A. Lowell, C. Sleator, J.A. Tomsick, A. Zoglauer, M. Amman, H.-K. Chang, C.-H. Tseng, C.-Y. Yang, C.-H. Lin, P. Jean, P. von Ballmoos, The 2016 Super Pressure Balloon flight of the Compton Spectrometer and Imager, arXiv:1701.05558, 2017.

\bibitem{Cumani2019LEO}
P. Cumani, M. Hernanz, J. Kiener, V. Tatischeff, A. Zoglauer,
Background for a gamma-ray satellite on a low-Earth orbit,
Exp. Astron. 47 (2019) 273--302,
doi:10.1007/s10686-019-09624-0.

\bibitem{Dobynde2021REDMoon}
M.I. Dobynde, J. Guo,
Radiation environment at the surface and subsurface of the Moon:
model development and validation,
J. Geophys. Res. Planets 126 (2021) e2021JE006930,
doi:10.1029/2021JE006930.

\bibitem{Takahashi2007HXD}
\revblue{T. Takahashi et al., Hard X-Ray Detector (HXD) on Board Suzaku,
Publ. Astron. Soc. Jpn. 59 (2007) S35--S51,
doi:10.1093/pasj/59.sp1.S35.}

\bibitem{Lotti2021Athena}
S. Lotti et al.,
Review of the particle background of the Athena X-IFU instrument,
Astrophys. J. 909 (2021) 111,
doi:10.3847/1538-4357/abd94c.

\bibitem{Verdier2011BGO}
\revblue{M.-A. Verdier, P.C.F. Di Stefano, P. Nadeau, C. Behan, M. Clavel, C. Dujardin,
Scintillation properties of $\mathrm{Bi_4Ge_3O_{12}}$ down to 3 K under
$\gamma$ rays,
Phys. Rev. B 84 (2011) 214306,
doi:10.1103/PhysRevB.84.214306.}

\bibitem{Usoskin2005}
I.G. Usoskin, K. Alanko-Huotari, G.A. Kovaltsov, K. Mursula,
Heliospheric modulation of cosmic rays: monthly reconstruction for 1951--2004,
J. Geophys. Res. Space Phys. 110 (2005) A12108,
doi:10.1029/2005JA011250.

\bibitem{Kuznetsov2017}
N.V. Kuznetsov, H. Popova, M.I. Panasyuk,
Empirical model of long-time variations of galactic cosmic ray particle fluxes,
J. Geophys. Res. Space Phys. 122 (2017) 1463--1472,
doi:10.1002/2016JA022920.

\bibitem{Cowan2011}
G. Cowan, K. Cranmer, E. Gross, O. Vitells, Asymptotic formulae for likelihood-based tests of new physics, Eur. Phys. J. C 71 (2011) 1554, doi:10.1140/epjc/s10052-011-1554-0.

\bibitem{Agnese2018NuclearRecoil}
R. Agnese et al., Nuclear-recoil energy scale in CDMS II silicon dark-matter detectors, Nucl. Instrum. Methods Phys. Res. A 905 (2018) 71--81, doi:10.1016/j.nima.2018.07.028.

\bibitem{Hull2023ThermalBackground}
S.V. Hull et al., Characterizing thermal background events for Athena X-IFU, IEEE Trans. Appl. Supercond. 33 (2023) 2100606, doi:10.1109/TASC.2023.3257131.

\end{thebibliography}
\end{document}
